\textsc{The aim of this chapter is to compare}
several notions of ``duality'' in monoidal categories.
On the one hand, we have closedness,
exemplifying the tensor–hom adjunction in the category of sets;
on the other lies rigidity,
generalising the duals of finite\hyp{}dimensional vector spaces.
\GV{} duality~\cite{boyarchenko13:groth-verdier},
also called \(*\)\hyp{}autonomy~\cite{barr79},
describes a duality theory between the strict confinements of rigidity
and the very general notion of monoidal closedness.
Roughly speaking, Grothendieck–Verdier categories
consist of a closed monoidal category and a chosen dualising object,
such that the internal hom into this object induces an anti-equivalence of categories.
While closed and rigid categories have left and right variants, \GV{} duality is an inherently ambidextrous concept.

Our first goal is
to more closely investigate the converse of \cref{prop: rigidity_yields_adjoints},
see also \cref{rmk:trep-is-not-rigid-first-exposure}.
That is, we study whether one can decide
if a closed monoidal category is rigid
solely by verifying that the internal hom is
given by tensoring with an object.
We thank Chris Heunen for suggesting this may not hold in general at the
\href{https://conferences.leeds.ac.uk/bcqt2022/}{\textsc{bcqt}\oldstylenums{2022}} summer school.

\begin{reptheorem}{thm:props-of-c}
  There exists a non-rigid tensor representable category.
\end{reptheorem}
\marginnote{%
  For the second inclusion
  we assume that the natural transformations ``comparing'' the left and right-handed version of tensor representability are invertible,
  see \cref{prop:trep-criteria,prop:right-tensor-rep->right-gv}.
  For rigid categories this is always the case.%
}
In particular,
taking
\cref{thm:props-of-c},
\cref{prop:right-tensor-rep->right-gv,ex:r-cat-not-always-trep},
as well as \cref{prop:gv->closed,rmk:closed-but-not-gv}
together yields a strict hierarchy:
\[
  \text{Rigid}
  \subsetneq \text{Tensor representable}
  \subsetneq \text{\GV}
  \subsetneq \text{Closed}.
\]

We then investigate closed and \GV{} structures on functor categories endowed with Day convolution as its tensor product.
Given a sufficiently nice base category, the functor category inherits these dualities.

\begin{reptheorem}{prop:fun-cat-GV}
  Let \(\cat{C}\) be a \(\Bbbk\)\hyp{}linear hom-finite Grothendieck–Verdier category.
  If \cref{hypothesis:tens-repr-mack} holds%
  —all relevant coends exist and are finite-dimensional vector spaces—%
  then \([\cat{C}, \kvect]\) is a Grothendieck–Verdier category.
\end{reptheorem}

We furthermore investigate the duality properties of Cauchy completions.
Since we are working with the functor category
\([\cat{C}, \kvect]\) instead of the presheaves \([\cat{C}^{\op}, \kvect]\),
we work with the Cauchy completion \(\overline{\cat{C}^{\op}}\) of \(\cat{C}^{\op}\),
it being a subcategory of the free cocompletion \(\Cpsh = [\cat{C}, \kvect]\) of \(\cat{C}^{\op}\).
Recall the definition of the Yoneda embedding from \cref{eq:yoneda-embedding}.

\begin{reptheorem}{cor:pres-ref-dual-struct}
  Let \(\cat{C}^{\op}\) be a \(\Bbbk\)\hyp{}linear right closed monoidal category.
  \begin{thmlist}[noitemsep]
    \item \(\cat{C}^{\op}\) is right rigid if and only if its Cauchy completion \(\overline{\cat{C}^{\op}}\) is.
    \item \(\cat{C}^{\op}\) is right tensor representable if and only if \(\overline{\cat{C}^{\op}}\) is.
    \item \((\cat{C}^{\op}, d)\) is a right \GV{} category if and only if \((\overline{\cat{C}^{\op}}, \yo_{d})\) is.
  \end{thmlist}
\end{reptheorem}

The chapter concludes with an investigation of three examples of abelian closed monoidal functor categories.
The first is finite Boolean algebras,
which historically arose in the study of order theory and universal algebra;
the connection to Grothendieck–Verdier duality was already noted in~\cite{barr79}.

\begin{reptheorem}{prop:qf2-algebras-from-boolean}
  The path algebra of a finite Boolean algebra is \textsc{qf}-2.%
  \marginnote{\normalfont{See \cref{def:qf-2} for when an algebra is \textsc{qf}-2.}}
\end{reptheorem}

We subsequently study the category of finite-dimensional Mackey functors,
objects abstracting and generalising the induction, restriction, and conjugation functors
in classical representation theory~\cite{lindner76:mackey,thevenaz95:mackey,webb00:mackey}.

\begin{reptheorem}{prop:mky-GV-and-rigid}
  Let \(G\) be a finite group and suppose \(\Bbbk\) is a field.
  The category \(\mkyfin_\Bbbk(G)\) is a \GV{} category,
  and its rigid objects are precisely the finitely\hyp{}generated projective Mackey functors.
  Furthermore, \(\mkyfin_{\Bbbk}(G)\) is rigid
  if and only if the Mackey algebra \(\mathbb{M}G\) is semisimple,
  which is equivalent to \(\mathrm{char}\,\Bbbk\) not dividing the order of \(G\).
\end{reptheorem}

Finally, we concern ourselves with strict 2-groups and their equivalent formulation as crossed modules~\cite{wagemann2021:CrossedModules}.
These are related to \emph{\(r\)-categories},
a special case of Grothendieck–Verdier duality
in which the dualising object is isomorphic to the monoidal unit.

\begin{reptheorem}{prop:2-groups-tensor-product}
  Let \(G\) and \(H\) be finite groups.
  Suppose that
  \[
    \big( G,\, H,\, t\from H \to G,\, \alpha \from G \to \mathrm{Aut}(H) \big)
  \]
  is a crossed module.
  Let \(\cat{G}\) be its associated strict \(2\)\hyp{}group.
  The category \(\mathsf{rep}_{\cat{G}}(\ker t)\)
  of finite-dimensional representations of \(\ker t\)
  is a right \(r\)\hyp{}category.
\end{reptheorem}

\section{Tensor representability}\label{sec:tensor-rep->gv}

\begin{definition}\label{def:tensor-representable-endo-functors}
  We call an endofunctor \(F\from \cat{C} \to\cat{C}\) of a monoidal category \emph{right tensor representable}
  if there exists an \(x \in\cat{C}\) such that \(F \cong x \otimes \blank \).
\end{definition}

\Cref{prop: rigidity_yields_adjoints} states that the internal hom of a rigid monoidal category is tensor-re\-pre\-sent\-able at every object.
From now on, for the sake of brevity, we will simply speak of ``(monoidal) tensor representable categories'' instead of ``closed monoidal categories with tensor representable internal homs''.
As shown in \cref{lem:condition-closed->tensor-dual} below, this is equivalent to the following definition, which does not rely on closedness.

\begin{definition}\label{def:tensor-representable}
  A monoidal category \(\cat{C}\) is called \emph{right tensor representable}
  if for all \(x \in \cat{C}\) there exists an object \(Rx \in\cat{C}\),
  called the \emph{right tensor dual} of \(x\),
  such that
  \(x \otimes \blank \adjoint Rx \otimes \blank\).
\end{definition}

A monoidal category \(\cat{C}\) is called \emph{left tensor representable} if \(\cat{C}^{\tensorop}\) is right tensor representable.
A left and right tensor repre\-sentable category will be referred to as \emph{tensor representable}.

Tensor representable categories encompass all rigid monoidal categories.
As such, some results in the theory of the latter carry over to the former.
For example, the following result is well-known in the case of rigid monoidal abelian categories,
see for example~\cite[Corollary~4.2.13]{Etingof2015}.

\begin{lemma}\label{lem:projectivity-in-tr-cats}
  All objects of an abelian monoidal right tensor representable category are projective
  if and only if its unit is projective.
\end{lemma}
\begin{proof}
  As the first claim implies the second, we only have to show its converse.
  Fix an object \(x\) in an abelian right tensor representable category \(\cat{A}\).
  There exists a chain of adjunctions
  \[
    x \otimes \blank \,\lad\, Rx \otimes \blank \,\lad\, R^2x \otimes \blank \,\lad\, \dots.
  \]
  Hence, the functor \(Rx \otimes \blank \from \cat{A} \to\cat{A}\) has left and right adjoints and is therefore exact.
  Consequently, \(\cat{A}(x, \blank) \cong \cat{A}(1, Rx \otimes \blank)\) is—as a composite of exact functors—exact itself.
\end{proof}

While \cref{lem:transfer-of-rigidity} holds in the rigid case,
it need not be true for tensor representable categories:
given a strong monoidal functor \(F\),
the image \(Fx\) of an object \(x\) admitting a (right) tensor dual may not have a (right) tensor dual itself.
However, the property is preserved by monoidal equivalences.

\begin{lemma}\label{lem:transfer-of-tensor-rep}
  Let \(\cat{C}\) be a right tensor representable category,
  and suppose that \(\cat{D}\) is monoidal.
  Any strong monoidal adjoint equivalence
  \(\adj{F}{F^{-1}}{\cat{C}}{\cat{D}}\)
  equips \(\cat{D}\) with a right tensor representable structure.
\end{lemma}
\begin{proof}
  For any \(x \in \cat{D}\), define \(Dx \defeq FRF^{-1}x \in\cat{D}\),
  where \(RF^{-1}x\) is a right tensor dual of \(F^{-1}x\in\cat{C}\).
  Then \(Dx \otimes \blank\) is right adjoint to \(x \otimes \blank\):
  \begin{align*}
    \cat{D}(x \otimes y, z)
    & \cong \cat{C}(F^{-1}(x \otimes y), F^{-1}z)
      \cong \cat{C}(F^{-1}x \otimes F^{-1}y, F^{-1}z) \\
    & \cong \cat{C}(F^{-1}y, RF^{-1}x \otimes F^{-1}z)
      \cong \cat{D}(FF^{-1}y, F(RF^{-1}x \otimes F^{-1}z)) \\
    & \cong \cat{D}(y, FRF^{-1}x \otimes z)
      = \cat{D}(y, Dx \otimes z).
  \end{align*}
  Thus, \(\cat{D}\) is right tensor representable.
\end{proof}

Just like in the rigid case,
the relationship between closedness and tensor repre\-sentability is governed by a family of natural isomorphisms.

\begin{lemma}\label{lem:condition-closed->tensor-dual}
  A monoidal category \(\cat{C}\) is right tensor representable
  if and only if
  it is right closed monoidal and for all \(x \in\cat{C}\) there are isomorphisms
  \begin{equation}\label{eq:closed->trep-merely-isos-required}
    \varphi^{(x)}_y \from{[x,y]}_r \to {[x,1]}_r \otimes y \qquad \qquad \text{natural in } y \in\cat{C}.
  \end{equation}
\end{lemma}
\begin{proof}
  For every \(x \in \cat{C}\) in a right tensor representable category \(\cat{C}\),
  by definition there exists an object \(Rx \in \cat{C}\) such that \(Rx\otimes \blank\) is right adjoint to \(x\otimes \blank\).
  Hence, the category \(\cat{C}\) is right closed monoidal,
  and the internal hom from \(x\in \cat{C}\) is given by \(Rx \otimes \blank\).
  The claim follows by noticing that \cref{eq:closed->trep-merely-isos-required}
  is now even the functorial equality  \(Rx \otimes \blank = Rx \otimes 1 \otimes \blank \).

  Conversely, for all \(x \in \cat{C}\) there is
  an isomorphism
  \[
    \cat{C}(y, {[x,z]}_r)
    \cong \cat{C}(x\otimes y, z)
    \cong \cat{C}(y, {[x,1]}_{r}\otimes z)
  \]
  natural in \(y,z \in \cat{C}\);
  the result follows by the Yoneda lemma.
\end{proof}

\begin{definition}\label{def:tensor-dual-functors}
  Assume \(\cat{C}\) to be a right tensor representable category.
  By \cref{lem:condition-closed->tensor-dual}, it is right closed monoidal.
  Then \(R \defeq {[\blank, 1]}_r \from \cat{C}^{\op} \to \cat{C}\) is called the \emph{right (tensor-)dualising functor} of \(\cat{C}\).
  Analogously, \(L \defeq {[ \blank , 1]}_\ell\from \cat{D}^{\op} \to \cat{D}\) is the \emph{left (tensor-)dualising functor} of a left tensor representable category \(\cat{D}\).
\end{definition}

A right rigid monoidal category is simultaneously left rigid if and only if its right dualising functor is an equivalence of categories.
This is more complicated in the tensor representable case.
Assume \(\cat{C}\) to be of this type.
Write \(R \from\cat{C}^{\op} \to\cat{C}\) and \(L \from\cat{C}^{\op} \to\cat{C}\) for the right and left tensor dualising functors, respectively.
Set
\begin{align*}
  \eta^{(x)}_y &\from y \to Rx \otimes x \otimes y, & \varepsilon^{(x)}_y&\from x \otimes Rx \otimes y \to y ,\\
  u^{(x)}_y & \from y \to y \otimes x \otimes Lx, & c^{(x)}_{y} &\from y \otimes Lx \otimes x \to y
\end{align*}
for the units and counits of the corresponding adjunctions.
The left and right tensor dualising functors are related to each other by
\begin{align}
  x \xrightarrow{u^{(Rx)}_x} x \otimes Rx \otimes LRx \xrightarrow{\varepsilon^{(x)}_1 \otimes LRx}LRx, \label{eq:can-morph-unit} \\
  x \xrightarrow{\eta^{(Lx)}_x} RLx \otimes Lx \otimes x \xrightarrow{RLx \otimes c^{(x)}_1}RLx. \label{eq:can-morph-counit}
\end{align}

\begin{proposition}\label{prop:trep-criteria}
  The right dualising functor \(R\from\cat{C}^{\op} \to\cat{C}\)
  of a tensor representable category \(\cat{C}\)
  is right adjoint to \(L^{\op}\from\cat{C}\to\cat{C}^{\op}\).
  Further, \(R\) is an equivalence if and only if
  the canonical morphisms of \cref{eq:can-morph-unit,eq:can-morph-counit} are invertible.
  In this case \(L^{\op}\) is a quasi-inverse of \(R\).
\end{proposition}
\begin{proof}
  For any two objects \(x,y \in\cat{C}\), we have
  \[
    \cat{C}^{\op}(Lx, y) =  \cat{C}(y, Lx) \cong  \cat{C} (y \otimes x, 1) \cong  \cat{C}(x,  Ry).
  \]
  A direct computation shows that the unit and counit of this adjunction are given by the natural transformations
  displayed in \cref{eq:can-morph-unit,eq:can-morph-counit}.
  If these morphisms are invertible, \(R\) and \(L^{\op}\) are quasi-inverse to each other.

  Conversely, assume \(R\) to be an equivalence of categories.
  Since any equivalence can be bettered to an adjoint equivalence
  and adjoints are unique up to unique natural isomorphism,
  \(L^{\op}\) must be quasi-inverse to \(R\)
  and the unit and counit morphisms of \cref{eq:can-morph-unit,eq:can-morph-counit} are invertible.
\end{proof}

\subsection{Grothendieck–Verdier duality}

\textsc{We are going to loosely follow}
the notation and terminology of~\cite{boyarchenko13:groth-verdier}.

\begin{definition}\label{def:*-autonomous-drinfeld}
  A \emph{(right) \GV{} category} consists of a pair \((\cat{C}, d)\) of a monoidal category \(\cat{C}\) and an
  object \(d\in \cat{C}\),
  such that for all \(x\in\cat{C}\) there exists a representing object \(D_r x \in \cat{C}\) for \(\cat{C}(x \otimes \blank, d)\);
  \ie,
  \begin{equation}\label{eq:*-aut-condition}
    \cat{C}(x \otimes \blank ,d) \cong \cat{C}(\blank, D_{r} x),
  \end{equation}
  and the induced \emph{right dualising functor} \(D_r\from \cat{C}^{\op}\to\cat{C}\) is an equivalence of categories.
  If \(d=1\) is the monoidal unit, one speaks of an \emph{\(r\)\hyp{}category}.
\end{definition}

Because \(D_r\) is an equivalence,
one also obtains
\begin{equation}\label{eq:gv-inverse-condition}
  \cat{C}(x \otimes y, d) \cong \cat{C}(y, D_r x) \cong \cat{C}^{\op}(D^{-1}_{\ell}y, x) = \cat{C}(x, D_r^{-1} y),
  \quad \text{for all } x, y \in \cat{C}.
\end{equation}

\begin{remark}\label{rmk:right-gv-vs-the-world}
  In~\cite{boyarchenko13:groth-verdier}, the authors instead consider
  a pair \((\cat{C}, d)\)
  subject to the natural isomorphism \(\cat{C}(\blank \otimes y, d) \cong \cat{C}(\blank, D_{\ell} y)\),
  such that the induced functor \(D_{\ell}\from \cat{C}^{\op} \to \cat{C}\) is an equivalence.
  In our framework, we shall call this a \emph{left} \GV{} category structure on \(\cat{C}\),
  with \emph{left} dualising functor \(D_{\ell}\).
  \Cref{eq:*-aut-condition,eq:gv-inverse-condition} become
  \begin{equation}\label{eq:left-gv-condition-and-inverse}
    \cat{C}(x \otimes y, d) \cong \cat{C}(x, D_{\ell}y) \cong \cat{C}(y, D^{-1}_{\ell}x).
  \end{equation}
\end{remark}

\begin{remark}\label{rmk:GV-is-symmetric}
  Notice that the nomenclature of \cref{rmk:right-gv-vs-the-world} might be misleading at first sight.
  Because the dualising functor of a \GV{} category is assumed to be an equivalence,
  \GV{} structures are always two-sided.
  In view of \cref{eq:gv-inverse-condition},
  a right \GV{} category \((\cat{C}, d)\)
  is also a left \GV{} category
  with the same dualising object \(d\), as well as \(D_r^{-1} = D_{\ell}\).
  Likewise, all left \GV{} categories are also right-sided.

  For the sake of clarity and to denote the directional bias,
  \GV{} structures will often nevertheless explicitly carry a prefix.
  In particular, in \cref{sec:tensor-rep-functor-cats}
  we consider the induced \GV{} structure of the category of copresheaves \(\Cpsh\) over a \GV{} base category \(\cat{C}\).
  In this case, the \emph{left} dualising functor of \(\cat{C}\) naturally gives rise to a \emph{right} dualising functor for \(\Cpsh\);
  see \cref{prop:fun-cat-GV}.
\end{remark}

The symmetric nature of \GV{} categories is in stark contrast to
the one\hyp{}sidedness of rigidity and tensor representability.

\begin{example}\label{ex:right-rigid-not-gv}
  Consider a Hopf algebra \(H\) whose antipode is not invertible,
  see~\cite{schauenburg2000:FaithfulFlatnessHopf}.
  Its finite\hyp{}dimensional right comodules are a left-rigid category \(\rcomod{H}\) which,
  as discussed in~\cite[Remark~2]{ulbrich1990:HopfAlgebrasRigid}, is not right rigid.
  If \(\rcomod{H}\) was a \GV{} category, there would
  as a consequence of \cref{rmk:gv-d-from-D}
  exist a comodule \(M \in \rcomod{H}\),
  such that \(\ld{\!(\blank)}\otimes M \from {(\rcomod{H})}^{\op}\to \rcomod{H}\) is an equivalence of categories.
  This would imply that \(M\) is one-dimensional,
  and therefore the left dualising functor \(\ld{\!(\blank)}\from \rcomod{H}^{\op}\to \rcomod{H}\) would be an equivalence—contradiction.
\end{example}

The following proposition, also discussed in~\cite[Section~2]{boyarchenko13:groth-verdier},
is a direct consequence of \cref{eq:*-aut-condition,eq:gv-inverse-condition}.

\begin{proposition}\label{prop:gv->closed}
  Every (right) \GV{} category \((\cat{C}, d)\) is closed mon\-oidal;
  for all \(x, y \in \cat{C}\), the left and right internal homs are given by
  \begin{equation}\label{eq:gv->internal hom}
    \begin{aligned}
      {[x, y]}_{\ell} &\defeq D_r^{-1}(x \otimes D_r y) = D_{\ell}(x \otimes D_{\ell}^{-1} y), \\
      {[x, y]}_r &\defeq D_r(D_r^{-1}y \otimes x) = D_{\ell}^{-1}(D_{\ell}y \otimes x).
    \end{aligned}
  \end{equation}
\end{proposition}

\begin{remark}\label{rmk:gv-d-from-D}
  If \((\cat{C},d)\) is a \GV{} category,
  one recovers its dualising object by
  applying the dualising functor to the monoidal unit of \(\cat{C}\).
  In fact, as shown in~\cite[Remark~2.1.(4)]{boyarchenko13:groth-verdier},
  this relationship is bidirectional:
  \begin{equation}\label{eq:gv:1:from:d:and:d:from:1}
    D_r 1 \cong d, \quad
    D_r d \cong 1, \quad
    D_{\ell} 1 \cong d, \quad
    D_{\ell} d \cong 1, \quad
    D_r^2 1 \cong 1, \quad
    D_{\ell}^2d \cong d.
  \end{equation}

  In combination with \cref{prop:gv->closed},
  one obtains an explicit description of the dualising functors of \(\cat{C}\):
  \begin{equation}\label{eq:gv-dualising-from-internal-hom-vice-versa}
    D_r x \cong {[x, d]}_r \qquad \text{ and } \qquad
    D_{\ell}x = D_r^{-1}x \cong {[x, d]}_{\ell}, \qquad \text{for all }x \in \cat{C}.
  \end{equation}
\end{remark}

While \GV{} duality is an extra bit of structure on a monoidal category%
—a priori making it quite distinct from the other dualities—%
in practice it still behaves very similarly to how a property would,
as the invertible objects govern the possible choices for dualising objects.

\begin{lemma}[{\cite[Proposition~2.3]{boyarchenko13:groth-verdier}}]\label{lem:choices-of-dualising-objects}
  Given a right \GV{} category \((\cat{C},d)\), there is an anti\hyp{}equiva\-lence between its full subcategories of invertible and dualising objects given by
  \[
    \mathrm{Inv}(\cat{C}) \to \mathrm{Dual}(\cat{C}), \qquad\alpha \mapsto \alpha^{-1}\otimes d.
  \]
\end{lemma}

\begin{example}\label{rmk:closed-but-not-gv}
  The category \((\mathsf{Set}, \times, \1)\) of sets
  with its usual Cartesian monoidal structure is closed monoidal.
  If some set \(X\) is invertible,
  then it must be isomorphic to the terminal set \(\1\).
  Thus, to deduce that there is no \GV{} structure on \(\Set\),
  it is sufficient to see that the functor \(\mathsf{Set}(\blank, \1)\) is not an anti\hyp{}equivalence,
  which follows by \(\1\) being terminal.
\end{example}

Let us now characterise the relationship between \GV{} and tensor representable categories.

\begin{proposition}\label{prop:right-tensor-rep->right-gv}
  Let \(\cat{C}\) be a tensor representable category.
  Then the following statements are equivalent:
  \begin{thmlist}[noitemsep]
    \item\label{itm: dualising-functor-trep-invertible} the right dualising functor admits a quasi-inverse,
    \item\label{itm: trep-is-r-cat} the category \(\cat{C}\) is an \(r\)\hyp{}category, and
    \item\label{itm: trep-is-GV} the category \(\cat{C}\) is a Grothendieck–Verdier category.
  \end{thmlist}
\end{proposition}
\begin{proof}
  By definition we have that~\ref{itm: trep-is-r-cat}\(\implies\)\ref{itm: trep-is-GV},
  and~\ref{itm: dualising-functor-trep-invertible}\(\implies\)\ref{itm: trep-is-r-cat} follows by \cref{def:tensor-dual-functors,rmk:gv-d-from-D}.

  To prove~\ref{itm: trep-is-GV}\(\implies\)\ref{itm: trep-is-r-cat},
  assume that \(\cat{C}\) is \GV{} with dualising object \(d\in \cat{C}\).
  Write \(L, R \from \cat{C}^{\op} \to \cat{C}\) for its left and right tensor dualising functors and
  \(D_r\from \cat{C}^{\op} \to \cat{C}\) for the dualising functor induced by \(d\in \cat{C}\).
  Then
  \begin{align*}
    Rd \otimes d
    \overset{\eqref{eq:closed->trep-merely-isos-required}}{\cong} {[d,d]}_r
    & \overset{\eqref{eq:gv->internal hom}}{\cong} D_r(D_r^{-1} d \otimes d)
      \overset{\eqref{eq:gv:1:from:d:and:d:from:1}}{\cong} D_r d
      \overset{\eqref{eq:gv:1:from:d:and:d:from:1}}{\cong} 1,\\
    d \otimes Ld
    \overset{\eqref{eq:closed->trep-merely-isos-required}}{\cong} {[d,d]}_{\ell}
    & \overset{\eqref{eq:gv->internal hom}}{\cong} D_r^{-1}(d \otimes D_r d)
      \overset{\eqref{eq:gv:1:from:d:and:d:from:1}}{\cong} D_r^{-1} d
      \overset{\eqref{eq:gv:1:from:d:and:d:from:1}}{\cong} 1,
  \end{align*}
  and therefore \(d \otimes Rd \cong d \otimes Rd \otimes d \otimes Ld \cong 1\).
  In other words, \(d\) is invertible and \(1 \cong Rd\otimes d\) is another dualising object due to \cref{lem:choices-of-dualising-objects}.

  To complete the proof notice that the right dualising functor of \(\cat{C}\) considered as an \(r\)\hyp{}category
  coincides with its right tensor dualising functor.
\end{proof}

\begin{example}\label{ex:r-cat-not-always-trep}
  The converse of \cref{prop:trep-criteria} is not true.
  Let \(A\) be a commutative Frobenius algebra over a field \(\Bbbk\).
  We write \(\lmod{A}^{\mathrm{f.g.}}\) for the category of finitely\hyp{}generated left \(A\)\hyp{}modules.
  One can show, see for example~\cite[Sections~15 and~16]{Lam1999}, that for all \(M \in \lmod{A}^{\mathrm{f.g.}}\) the canonical morphism
  \[
    \phi_M \from M \to \rrd{M} = \Hom_A(\Hom_A(M,A),A) , \qquad \qquad \phi_M(x)\alpha = \alpha(x)
  \]
  is an isomorphism;
  the module \(M\) is called \emph{reflexive}.

  Put differently, \(\lmod{A}^{\mathrm{f.g.}}\) is an \(r\)\hyp{}category with
  \[
    \Hom_A(\blank,A) \from {(\lmod{A}^{\mathrm{f.g.}})}^{\op} \to \lmod{A}^{\mathrm{f.g.}}
  \]
  as dualising functor.
  It is well-known that
  for any \(A\)\hyp{}module \(M\)
  there exists an \(A\)\hyp{}module \(N\)
  such that \(M \otimes \blank \lad N \otimes \blank\) if and only if \(M\) is finitely\hyp{}generated projective;
  see for example~\cite[Proposition~2.1]{niefield17:coexp}.
  Thus, \(\lmod{A}^{\mathrm{f.g.}}\) is tensor representable if and only if
  all (finitely\hyp{}generated) \(A\)\hyp{}modules are projective,
  which is equivalent to \(A\) being semisimple.

  However, not all commutative Frobenius algebras are semisimple.
  For example, the group algebra \(\Bbbk[G]\) of a finite abelian group \(G\) is Frobenius.
  By Maschke's theorem, it is semisimple if and only if the characteristic of \(\Bbbk\) does not divide the order of \(G\).
\end{example}

\begin{remark}
  Let \(\cat{C}\) be a right tensor representable category whose dualising functor \(R \from\cat{C}^{\op} \to\cat{C}\) is an equivalence.
  Following~\cite[Section~4]{boyarchenko13:groth-verdier},
  for all \(x,y \in\cat{C}\) there exists a canonical isomorphism
  \begin{align*}
    \tau_{x,y}\from \cat{C}(y\otimes x\otimes Rx \otimes Ry, 1)
    &\iso \cat{C}(Rx \otimes Ry, R(y \otimes x)) \\
    &\iso \cat{C}^{\op}(y \otimes x, R^{-1}(Rx \otimes Ry)).
  \end{align*}

  As shown in~\cite[Lemma~4.1 and Remark~4.2]{boyarchenko13:groth-verdier},
  the morphisms
  \[
    f \defeq \tau_{x,y}\big(\varepsilon^{(y)}_1 \circ (y \otimes \varepsilon^{(x)}_1 \otimes Ry)\big)\from R^{-1}(Rx \otimes Ry) \to y \otimes x
  \]
  \[
    \mu_{x,y} \defeq Rf \from Rx \otimes Ry \to R(y \otimes x) = R(x \otimes^{\op} y)
  \]
  endow \(R\from\cat{C}^{\op,\tensorop} \to\cat{C}\) with the structure of a lax monoidal functor:
  \begin{align*}
    &\mu_{x, y \otimes^{\op} z} \circ (Rx \otimes \mu_{y,z}) = \mu_{x \otimes^{\op} y, z} \circ (\mu_{x,y}\otimes Rz), \\
    &\mu_{1,x} \circ (R1 \otimes Rx) = \id_{Rx} = \mu_{x,1} \circ (Rx \otimes R1),
      \qquad\text{for all \(x,y,z \in\cat{C}\)}.
  \end{align*}
\end{remark}

In the situation of the above remark and
as a consequence of~\cite[Proposition~4.4]{boyarchenko13:groth-verdier},
we obtain the following characterisation of rigidity.

\begin{proposition}\label{prop:GV-rigid-characterisation}
  Let the dualising functor \(R \from\cat{C}^{\op, \tensorop} \to\cat{C}\) of a right tensor representable category be an equivalence.
  Then \(\cat{C}\) is rigid if and only if the morphism
  \(\mu_{x,y} \from R(\blank) \otimes R(\bblank) \nt R(\blank \otimes^{\op} \bblank)\) is an isomorphism,
  making \(R\) strong monoidal.
\end{proposition}

\subsection{The free tensor representable category is not rigid}\label{sec:negative}

\textsc{The aim of this section is} to concretise \cref{rmk:trep-is-not-rigid-first-exposure}.
That is, we explicitly construct a ``free'' tensor representable category \(\cat{W}\) and show that it is not rigid;
see \cref{thm:props-of-c}.
Due to the reliance of the proof on some technical arguments,
we will first provide an overview of the underlying strategy.

As an intermediate step in the construction of \(\cat{W}\), we define a strict monoidal category \(\cat{V}\), which has the same generators as \(\cat{W}\) but fewer relations.
In other words, \(\cat{W}\) itself will be a quotient of \(\cat{V}\).
The category \(\cat{V}\) will have the property that all presentations of a given morphism have the same length.
In order to show that \(\cat{W}\) is not rigid, we will fix an object \(t \in\cat{W}\) and assume it admits a rigid dual.
In particular, \(t\) would have to admit evaluation and coevaluation morphisms \(\ev_t \from t \otimes \rd t \to 1\) and \( \coev_t \from  1 \to \rd{t} \otimes t\).
By showing that any morphism in \(\cat{V}\) representing the class
\((\ev_t \otimes t) \circ (t \otimes \coev_{t}) \in \cat{W}(t,t)\) has to have length at least 2, we conclude that the snake identities cannot hold.
Hence, \(\cat{W}\) is not rigid.

We define the strict monoidal category \(\cat{V}\) in terms of generators and relations.
For details of this type of construction we refer to~\cite[Chapter~\textsc{xii}]{Kassel1998}.

The monoid of objects of the category \(\cat{V}\) is freely generated by the alphabet
\(\{\, \dots, R^{-1}t,\, t,\, Rt, \dots \,\}\).
That is, its elements are words of finite length whose letters are \(R^n t\) for some integer \(n \in\mathbb{Z}\).
Write \(1\in\Ob(\cat{V})\) for the monoidal unit of \(\cat{V}\), given by the empty word.
For an object \(w \defeq (R^{k_1}t, \dots, R^{k_n}t) \in \Ob(\cat{V})\), set \(Rw \defeq (R^{k_n+1}t, \dots, R^{k_1+1}t)\in \Ob(\cat{V})\).
This assignment extends to a group action of the free abelian group \(\langle R \rangle \cong (\mathbb{Z},+)\) on \(\Ob(\cat{V})\).
For any \(i \in\mathbb{Z}\), we write \(R^{i}w\) for the action of \(R^{i}\) on \(w \in \Ob(\cat{V})\).

The morphisms of \(\cat{V}\) are tensor products and compositions of identities and the \emph{generating morphisms}.
These are given for all \(x \neq 1, y \in \Ob(\cat{V})\) by
\begin{align*}
  \eta^{(x)}_{y}&\from y \to Rx \otimes x \otimes y,
  & \varepsilon^{(x)}_{y} &\from x \otimes Rx \otimes y\to y,\\
  u^{(x)}_{y} &\from y \to y \otimes x \otimes R^{-1}x,
  & c^{(x)}_{y} &\from y\otimes R^{-1}x \otimes x \to y.
\end{align*}
These relations are tailored to implement natural transformations
\begin{align*}
  \eta^{(x)} &\from \Id_{\cat{V}} \nt Rx \otimes x \otimes \blank, &
  \varepsilon^{(x)} &\from x \otimes Rx \otimes \blank \nt \Id_{\cat{V}}, \\
  u^{(x)} &\from \Id_{\cat{V}} \nt \blank \otimes x \otimes R^{-1}x, &
  c^{(x)} & \from \blank \otimes R^{-1}x \otimes x \nt \Id_{\cat{V}}.
\end{align*}
Explicitly,
for any generating morphism \(g \from c\to d\),
any \(y = a \otimes c \otimes b\), \(z= a \otimes d \otimes b\),
any arrow \(f = (\id_a \otimes g \otimes \id_b) \from y\to z\),
and all \(x \in \Ob(\cat{V}) \setminus \{1\}\),
we require the following conditions to hold:
\begin{align*}
  \eta^{(x)}_{z} \circ f &= (Rx \otimes x \otimes f) \circ \eta^{(x)}_{y},
  & \varepsilon^{(x)}_{z} \circ (x \otimes Rx \otimes f) &= f \circ \varepsilon^{(x)}_{y},\\
  u^{(x)}_{z} \circ f &= (f \otimes x \otimes R^{-1}x) \circ u^{(x)}_{y},
  & c^{(x)}_{z} \circ (f \otimes R^{-1} x\otimes x) &= f \circ c^{(x)}_{y}.
\end{align*}

We obtain the tensor representable category \(\cat{W}\)
by quotienting out the triangle identities.\note{%
  The category is obtained by \emph{monoidal localisation}, see~\cite{day73:note}.%
}
Explained in more detail, the strict monoidal category \(\cat{W}\) has the same objects and generating morphisms as \(\cat{V}\),
and the same identities hold.
In addition, for any \(x,y \in \Ob(\cat{W})\) with \(x \neq 1\), we require
\begin{equation}\label{eq:defining-relations-of-c}
  \begin{aligned}
    \varepsilon^{(x)}_{x \otimes y} \circ (x \otimes \eta^{(x)}_y) &= \id_x, &
    (Rx \otimes \varepsilon^{(x)}_{y}) \circ \eta^{(x)}_{Rx \otimes y} &= \id_{Rx},\\
    c^{(x)}_{y\otimes x} \circ (u^{(x)}_{y} \otimes x) &= \id_{x}, &
    (c^{(x)}_{y} \otimes R^{-1}x) \circ u^{(x)}_{y \otimes R^{-1}x}&= \id_{R^{-1}x},\\
    (\varepsilon^{(x)}_1 \otimes x) \circ u^{(Rx)}_{x} &= \id_x, &
    (x \otimes c^{(x)}_1) \circ \eta^{(R^{-1}x)}_{x} &= \id_x.
  \end{aligned}
\end{equation}

The next result succinctly summarises the observations made so far concerning the internal hom of \(\cat{W}\).

\begin{lemma}\label{lemma:free-trep-cat}
  The category \(\cat{W}\) is tensor representable.
  Furthermore, its right dualising functor \(R\from \cat{W}^{\op}\to\cat{W}\) is an equivalence of categories.
\end{lemma}
\begin{proof}
  By construction, we have \(x \otimes \blank \lad Rx \otimes \blank\) and \(\blank \otimes x \lad \blank \otimes R^{-1}x\) for all \(x \in\cat{W}\).
  Thus, \(\cat{W}\) is tensor representable.
  \Cref{eq:defining-relations-of-c,prop:trep-criteria}
  imply that \(R \from\cat{W}^{\op} \to\cat{W}\) and \(R^{-1} \from\cat{W}^{\op} \to\cat{W}\) are quasi-inverses.
\end{proof}

To analyse the morphisms in \(\cat{W}\) and show that it is not rigid monoidal, we will rely on two tools.
The first is the \emph{length} of an arrow \(f\in \cat{W}(x,y)\).
It is defined as the minimal number of generating morphisms needed to present \(f\).
The second tool will be given by invariants for morphisms in \(\cat{W}\) arising from functors into the rigid monoidal category \(\kvect\) of finite\hyp{}dimensional vector spaces over a field \(\Bbbk\).
We will write \(\rd*{(\blank)} \from \kvect^{\op} \to \kvect\) for the right rigid dualising functor of \(\kvect\).
The \(n\)\hyp{}fold dual of a finite\hyp{}dimensional vector space \(V\) will be denoted by \(V^{(n)}\).

\begin{lemma}\label{lem:kassel-tangle-strong-mon-fun}
  For all \(V\in\kvect\) there is a strong monoidal functor \(F_V \from\cat{W} \to \kvect\) satisfying
  \(F_V(R^{n}t) \cong V^{(n)}\) for all \(n \in \mathbb{Z}\) and
  \begin{align*}
    F_{V}(\eta^{(x)}_{y}) &= \coev^r_{F_{V}x} \otimes F_{V}y,
    & F_{V}(\varepsilon^{(x)}_{y}) &=\ev^r_{F_{V}x}\otimes F_{V}y,\\
    F_{V}(u^{(x)}_{y}) &= F_{V}y \otimes \coev^{\ell}_{F_{V}x},
    & F_{V}(c^{(x)}_{y}) &= F_{V}y \otimes \ev^{\ell}_{F_{V}x},
  \end{align*}
  for suitably chosen evaluation and coevaluation morphisms.
\end{lemma}
\begin{proof}
  \Cref{thm: rigid_strictification} yields a monoidal equivalence \(\adj{G}{H}{\kvect}{\cat{A}}\)
  between the category of finite\hyp{}dimensional vector spaces and a rigid monoidal category \(\cat{A}\)
  whose dualising functors \(D_{\ell}, D_r \from \cat{A}^{\op}\to \cat{A}\) satisfy
  \(D_{\ell}D_r^{\op}=\Id_{\cat{A}}=D_r^{\op}D_{\ell}\).
  Fix a quasi-inverse \(H \from \cat{A} \to \kvect\) and choose an object \(v\in \cat{A}\) with \(H(v) \cong V\), for some \(V \in \kvect\).
  A direct computation using the universal property of \(\cat{W}\), see for example~\cite[Proposition~\textsc{xii}.1.4]{Kassel1998}, shows that
  for any object \(v\in \cat{A}\) there is a unique strict monoidal functor \(F'_{v}\from \cat{W}\to \cat{A}\) such that \(F'_{v}(t)=v\) and for all \(x,y\in \cat{W}\), \(x\neq 1\), we have
  \begin{align*}
    F'_{v}\eta^{(x)}_{y} &= \coev^r_{F'_{v}x} \otimes F'_{v}y, &
    F'_{v}\varepsilon^{(x)}_{y} &=\ev^r_{F'_{v}x}\otimes F'_{v}y,  \\
    F'_{v}u^{(x)}_{y} &= F'_{v}y \otimes \coev^{\ell}_{F'_{v}x}, &
    F'_{v}c^{(x)}_{y} &= F'_{v}y \otimes \ev^{\ell}_{F'_{v}x}.
  \end{align*}
  The claim now follows by setting \(F=HF'_{G(V)}\).
\end{proof}

\begin{corollary}\label{cor: structure-of-isos}
  The following arrows cannot be isomorphisms,
  for any \(1 \neq x \in\cat{W}\) and any morphism \(g \in \cat{W}(a_1 \otimes y \otimes b_1,a_2 \otimes z \otimes b_2)\):
  \[
    (\id_{a_2} \otimes u^{(x)}_{z} \otimes \id_{b_2}) \circ g,
    \qquad \qquad
    g \circ (\id_{a_1} \otimes \varepsilon^{(x)}_{y} \otimes \id_{b_1}).
  \]
\end{corollary}
\begin{proof}
  Suppose that \(g \in \cat{W}(a_1 \otimes y \otimes b_1,a_2 \otimes z \otimes b_2)\)
  and consider
  \[
    f \defeq g \circ (\id_{a_1} \otimes \varepsilon^{(x)}_{y} \otimes \id_{b_1}).
  \]
  Let \(V\in\kvect\) be a vector space of dimension at least two.
  Applying \(F_V\) to \(f\), we get
  \(F_{V}f = F_{V}g \circ F_{V}(\id_{a_1} \otimes \varepsilon^{(x)}_{y} \otimes \id_{b_1})\).
  However, due to the difference in the dimensions of its source and target,
  \(F_{V}(\id_{a_1} \otimes \varepsilon^{(x)}_{y} \otimes \id_{b_1})\) must have a non-trivial kernel and thus \(f\) cannot be an isomorphism.

  A similar argument involving the cokernel proves that the composition
  \((\id_{a_2} \otimes u^{(x)}_{z} \otimes \id_{b_2}) \circ g\)
  is not invertible.
\end{proof}

\begin{theorem}\label{thm:props-of-c}
  The category \(\cat{W}\) is not rigid.
\end{theorem}
\begin{proof}
  Assume that \(t\in\cat{W}\) admits a right rigid dual \(x\in\cat{W}\).
  By the uniqueness of adjoints,
  there exist isomorphisms \(\vartheta \from Rt \otimes t \to x \otimes t\) and \(\theta \from x \to Rt\) such that the evaluation and coevaluation morphisms are given by
  \[
    \coev \defeq \vartheta \circ \eta^{(t)}_{1} \from 1 \to x \otimes t, \qquad \qquad
    \ev \defeq \varepsilon^{(t)}_{1} \circ (t\otimes \theta) \from t \otimes x \to 1.
  \]

  Let \(\pi \from  \cat{V} \sur \cat{W}\) be the projection functor.
  We consider the set
  \[
    S \defeq \left\{
      (\varepsilon^{(t)}_1 \otimes t)\,\phi\,(t \otimes \eta^{(t)}_1) \in \cat{V}(t,t) \,\middle\vert\, \begin{gathered}\phi \in \cat{V}(t\otimes Rt \otimes t, t \otimes Rt \otimes t) \\\text{ such that }  \pi(\phi) \text{ is invertible} \end{gathered} \right\}.
  \]
  Notice that \(S \subseteq \Mor(\cat{V})\).
  By construction, there exists a morphism \(s\in S\) such that the composite arrow
  \((\ev \otimes t) \circ (t\otimes \coev) = \pi(s)\)
  corresponds to one of the snake identities.
  Furthermore, every element of \(S\) has length at least two.\note{%
    Note that the relations of \(\cat{V}\) leave the number of generating morphisms in any presentation of a given arrow invariant.
  }
  Thus, by proving that \(S\) is closed under the relations arising from \cref{eq:defining-relations-of-c},
  it follows that \(\pi(s)\neq \id_t\), which concludes the proof.

  Let us consider an arbitrary element \(f \defeq (\varepsilon^{(t)}_1 \otimes t) \circ \phi \circ (t \otimes \eta^{(t)}_1) \in S\).
  There are two types of ``moves'' we have to study.
  First, suppose we expand an identity into one of the morphisms displayed in \cref{eq:defining-relations-of-c}.
  This equates to either pre or postcomposing \(\phi\) with an arrow \(\psi \in \cat{V}(t \otimes Rt \otimes t,t \otimes Rt \otimes t)\)
  that projects onto an isomorphism in \(\cat{W}\), leading to another element in \(S\).

  Second, any of the morphisms of \cref{eq:defining-relations-of-c} might be contracted to an identity.
  A priori, there are three ways in which this might occur:
  \begin{align}
    f&= (\varepsilon^{(t)}_1 \otimes t) \circ \phi' \circ \varepsilon^{(t)}_t \circ (t \otimes \eta^{(t)}_1),  &&\!\!\!\text{with } \phi=\phi' \circ \varepsilon^{(t)}_t \text{, or}\label{eq:bad-case-1}\\
    f&= (\varepsilon^{(t)}_1 \otimes t) \circ u^{(Rt)}_t \circ \phi'' \circ (t \otimes \eta^{(t)}_1), &&\!\!\!\text{with } \phi = u^{(Rt)}_t \circ \phi'' \text{, or}\label{eq:bad-case-2}\\
    f&= (\varepsilon^{(t)}_1 \otimes t) \circ \phi_2 \circ g \circ \phi_1 \circ (t \otimes \eta^{(t)}_1), &&\!\!\!\text{with } \phi=\phi_2 \circ g \circ \phi_1 \text{ and \(\pi(g)=\id\).}
  \end{align}
  Due to \cref{cor: structure-of-isos},
  neither \(\pi(\phi') \circ \pi(\varepsilon^{(t)}_{t})\) nor \(\pi(u^{(Rt)}_{t}) \circ \pi(\phi'')\) are isomorphisms,
  contradicting Cases~\eqref{eq:bad-case-1} and~\eqref{eq:bad-case-2}.

  Now, assume \(f = (\varepsilon^{(t)}_1\otimes t) \circ \phi_2 \circ g \circ \phi_1 \circ (t \otimes \eta^{(t)}_1)\) and \(\phi=\phi_2 \circ g \circ \phi_1\).
  Using the functoriality of \(\pi \from \cat{V} \to \cat{W}\), we get
  \[
    \pi(\phi)
    = \pi(\phi_2 \circ g \circ \phi_1)
    = \pi(\phi_2) \circ \pi(g) \circ \pi(\phi_1)
    = \pi(\phi_2) \circ \pi(\phi_1)
    = \pi(\phi_2 \circ \phi_1).
  \]
  Thus, \(\pi(\phi_2 \circ \phi_1)\) is an isomorphism
  and \((\varepsilon^{(t)}_1\otimes t) \circ \phi_2 \circ \phi_1 \circ (t \otimes \eta^{(t)}_1)\) is an element of \(S\).
\end{proof}

\section{Functor categories}\label{sec:tensor-rep-functor-cats}

\textsc{We will now investigate} the previously discussed types of dualities
in the context of functor categories.
The starting point is the fact that Day convolution is closed monoidal,
see \cref{thm:closed-mon-str-on-functors}.
Examples of such categories of functors
include modules over path algebras, Mackey functors, and
group-graded representations arising from crossed modules;
we shall more closely investigate these in \cref{sec:mky-2-grps}.

\begin{hypothesis}\label{hypo:eso-enriched}
  Fix a field \(\Bbbk\),
  and write \(\rd*{(\blank)}\from \kvect^{\op} \to \kvect\) for the \(\Bbbk\)\hyp{}linear dualising functor.
  Unless explicitly stated,
  for the remainder of this section
  all (monoidal) categories and functors will be enriched over \(\kVect\),
  and \(\cat{C}\) will denote an essentially small monoidal \(\Bbbk\)\hyp{}linear category.
  For the sake of brevity, the prefixes ``enriched'', ``\(\Bbbk\)\hyp{}linear'', and ``(essentially) small'' will often be omitted.
  The category of functors from \(\cat{C}\) to \(\kVect\) will be denoted by \(\Cpsh\).
\end{hypothesis}

Instead of studying functors \(F \in \Cpsh = [\cat{C}, \kVect]\)
from a hom-finite category \(\cat{C}\) to \emph{all} vector spaces,
we are interested in the full subcategory \(\Cpsh_{\fin}\) of (point-wise)
\emph{finite\hyp{}dimensional} functors.
That is, \(F \in \Cpsh\) such that \(Fx\) is finite-dimen\-sional for all \(x \in \cat{C}\).

\begin{hypothesis}\label{hypothesis:tens-repr-mack}
  Within this section, we assume that all relevant ends and coends are finite\hyp{}dimensional vector spaces.
\end{hypothesis}

In particular, \cref{hypothesis:tens-repr-mack}
implies that \(\Cpsh_{\fin}\) is closed monoidal;
that is, the tensor product and the internal hom are finite\hyp{}dimensional at every point.
These assumptions impose a certain finiteness condition on \(\cat{C}\) itself, or—as the next example will show—on a dense\note{%
  A subcategory \(\cat{O}\) of \(\cat{C}\) is \emph{dense} if
  the restricted Yoneda embedding from \(\cat{C}\) to \([\cat{O}^{\op}, \kVect]\)
  is fully faithful.%
}
subcategory of it.

\begin{example}\label{ex:mackey-dense-subcategories}
  Let \(\cat{C}\) be a hom-finite category,
  and assume that there is a full subcategory \(\cat{S}\) such that every object of \(\cat{C}\) may be written as a finite direct sum of objects in \(\cat{S}\).
  In this setting, one can reformulate \cref{lem:enriched-yoneda} to index over this finite set:
  for example, given \(F\from \cat{C}\to \kvect\), we have
  \begin{align*}
    Fx &\cong F\Big(\!\bigoplus_{i = 1}^n s_i\Big)
         \cong \bigoplus_{i = 1}^n Fs_i
         \cong \bigoplus_{i = 1}^n \int^{s \in \cat{S}}\!\! \cat{S}(s, s_i) \otimes Fs
         \cong \bigoplus_{i = 1}^n \int^{s \in \cat{S}}\!\! \cat{C}(s, s_i) \otimes Fs \\
       &\cong \int^{s \in \cat{S}}\!\! \cat{C}\Big(s, \bigoplus_{i = 1}^n s_i\Big) \otimes Fs
         \cong \int^{s \in \cat{S}}\!\! \cat{C}(s, x) \otimes Fs.
  \end{align*}
  Panchadcharam and Street used this to show that finite\hyp{}dimensional Mackey functors are a \GV{} category, see~\cite[Section~9]{panchadcharam07:mackey}.
\end{example}

Hypothesis~\ref{hypothesis:tens-repr-mack}
does not only imply a closed structure on \(\Cpsh_{\fin}\),
but yields a canonical notion of a dual.

\begin{proposition}\label{prop:fun-cat-GV}
  Let \((\cat{C}, d)\) be a hom-finite left \GV{} category with dualising functor \(D_{\ell}\).
  Then \(\Cpsh_{\fin}\) is a right \GV{} category with dualising object
  \(\rd*{\cat{C}(\blank, d)} \in \Cpsh_{\fin}\) and dualising functor
  \[
    \mathfrak{D}_r \from \Cpsh_{\fin}^{\op} \to \Cpsh_{\fin}, \qquad \qquad
    F \mapsto {F(D_{\ell}\blank)}^{*}.
  \]
\end{proposition}
\begin{proof}
  Using \cref{eq:gv-dualising-from-internal-hom-vice-versa},
  one can recover the right dualising functor of a \GV{} category from its right internal hom
  by evaluating it at the dualising object;
  \ie, \(\mathfrak{D}_r \cong {[\blank, {\cat{C}(\blank, d)}^{*}]}_r\).
  In our case, we have
  \begin{align*}
    {[F, \rd*{\cat{C}(\blank, d)}]}_r x
    \overset{\eqref{eq:right-internal hom-simplification}}
    &{\cong} \int_b \kvect\big(F{[x,b]}_{\ell}, \kvect(\cat{C}(b,d), \Bbbk)\big) \\
    & \cong \int_b \kvect\big(\cat{C}(b, d) \otimes F{[x,b]}_{\ell}, \Bbbk\big) \\
    & \cong \kvect\Big(\!\int^b\!\!\cat{C}(b,d)\otimes F{[x,b]}_{\ell}, \Bbbk \Big)
      \overset{\eqref{eq:coYoneda-covariant}}{\cong} \kvect\big(F{[x,d]}_{\ell}, \Bbbk\big)\\
    \overset{\eqref{eq:gv-dualising-from-internal-hom-vice-versa}}
    &{\cong} \rd*{F(D_{\ell}x)}.
  \end{align*}
  To conclude the proof, notice that  \(\mathfrak{D}_r \from \Cpsh_{\fin}^{\op}\to \Cpsh_{\fin}\) has a quasi-inverse
  \[
    \mathfrak{D}_r^{-1} \from\Cpsh_{\fin} \to \Cpsh_{\fin}^{\op}, \qquad \qquad F \mapsto \rd*{F(D_{\ell}^{-1}\blank)}.\qedhere
  \]
\end{proof}

While we will develop efficient means to detect tensor representability and rigidity based on properties which exist in the abelian case,
it is worthwhile to explore how such structures arise as an interplay between the duality type of the base category and that of \(\kvect\).

\begin{corollary}\label{cor:day-compact}
  \marginnote{\textnormal{%
    Not all assumptions are strictly needed, see \cref{rmk:reduction-of-assumptions}.}%
  }
  Let \((\cat{C}, d)\) be a hom-finite left \GV{} category
  with dualising functor \(D_{\ell}\).
  Then the category \(\Cpsh_{\fin}\) is right tensor representable if there are isomorphisms
  \begin{equation}\label{eq:trep-nice-enough}
    D_{\ell}^2a \cong a,
    \qquad D_{\ell}(a\otimes x) \cong D_{\ell}x \otimes D_{\ell}a,
    \qquad \text{ natural in } a,x \in \cat{C},
  \end{equation}
  and for all objects \(F, G \in \Cpsh_{\fin}\) we have
  \begin{equation}\label{eq:end-coend-commute}
    \int^a \kvect(F{[\blank,a]}_{\ell}, Ga) \cong \int_a \kvect(F{[\blank,a]}_{\ell},Ga).
  \end{equation}
\end{corollary}
\begin{proof}
  Let \(\mathfrak{D}_r\) be the dualising functor of \cref{prop:fun-cat-GV},
  and fix functors \(F, G \in \Cpsh_{\fin}\).
  Then, by the following calculation,
  we obtain an isomorphism
  \((\mathfrak{D}_r F * G) x \cong {[F,G]}_r x\),
  and hence the claim follows:
  \begin{align*}
    (\mathfrak{D}_r F& * G) x
    \overset{\eqref{eq:day-convolution-simp-left}}{\cong}
      \int^a\mathfrak{D}_r F({[a,x]}_{\ell}) \otimes Ga
      \overset{\eqref{eq:gv->internal hom}}{\cong}
      \int^a\mathfrak{D}_r F(D_{\ell}(a \otimes D_{\ell}x)) \otimes Ga \\
    \overset{\ref{prop:fun-cat-GV}}%
    &{\cong}
      \int^a\rd*{F(D_{\ell}^2(a \otimes D_{\ell}x))} \otimes Ga
      \overset{\eqref{eq:trep-nice-enough}}{\cong}
      \int^a\kvect\big(F(a \otimes D_{\ell}x), Ga\big) \\
    \overset{\eqref{eq:trep-nice-enough}}%
    &{\cong}
      \int^a\kvect\big(F(D^2_{\ell} a \otimes D_{\ell} x), Ga\big)
      \overset{\eqref{eq:trep-nice-enough}}{\cong}
      \int^a\kvect\big(F(D_{\ell}(x \otimes D_{\ell} a)), Ga\big) \\
    \overset{\eqref{eq:gv->internal hom}\,+\,\eqref{eq:trep-nice-enough}}%
    &{\cong}
      \int^a\kvect\big(F({[x,a]}_{\ell}), Ga\big)
      \overset{\eqref{eq:end-coend-commute}}{\cong}
      \int_a\kvect\big(F({[x,a]}_{\ell}), Ga\big) \\
    \overset{\eqref{eq:left-internal hom-simplification}}%
    &{=} {[F,G]}_r x.\qedhere
  \end{align*}
\end{proof}

\begin{remark}\label{rmk:conditions-for-early-properties}
  Before we discuss conditions for the interchangeability of ends and coends,
  let us briefly mention some cases where the dualising functor of a left \GV{} category \((\cat{C}, d)\)
  admits natural isomorphisms as stated in \cref{eq:trep-nice-enough}.

  The requirement \(D_{\ell}(x \otimes y) \cong D_{\ell}y \otimes D_{\ell}x\) implies that \(\cat{C}\) must be left tensor repre\-sentable,
  since then for all \(x,y, z \in\cat{C}\) we have
  \[
    \cat{C}(x \otimes y, z)
    \overset{\eqref{eq:left-gv-condition-and-inverse}}{\cong}
    \cat{C}(x \otimes y \otimes D_{\ell}^{-1}z, d)
    \overset{\eqref{eq:left-gv-condition-and-inverse}}{\cong}
    \cat{C}(x, D_{\ell}(y \otimes D_{\ell}^{-1}z))
    \overset{\eqref{eq:trep-nice-enough}}{\cong}
    \cat{C}(x,z \otimes  D_{\ell}y).
  \]
  In the absence of further coherence assumptions,
  this does not imply that \(\cat{C}\) must be rigid.
  The condition \(D_{\ell}^2 \cong \Id_{\cat{C}}\) is met for example if \(\cat{C}\) is braided.
  For all \(x, y \in\cat{C}\) we have
  \[
    \cat{C}(x, y)
    \overset{\eqref{eq:left-gv-condition-and-inverse}}{\cong}
    \cat{C}(D_{\ell}y \otimes x,d)
    \cong
    \cat{C}(x \otimes D_{\ell}y  ,d)
    \overset{\eqref{eq:left-gv-condition-and-inverse}}{\cong}
    \cat{C}(x, D_{\ell}^2 y),
  \]
  and the Yoneda lemma implies that \(\Id_{\cat{C}} \cong D_{\ell}^2\).
\end{remark}

The following constructions are an adaptation of~\cite{day06:compac}.
For simplicity, we replaced the promonoidal structure in \emph{ibid}
with the one induced from a monoidal structure.
That is, \(J \defeq \cat{C}(1, \blank )\), \(P(x, y, z) \defeq \cat{C}(x \otimes y, z)\),
for \(J\) and \(P\) the pro\-mon\-oidal unit and multiplication.

\begin{remark}\label{rmk:preconditions-day-yoneda}
  Let \(\cat{C}\) be a category.
  Given any functor \(T\from \cat{C}^{\op} \otimes_{\Bbbk} \cat{C} \to \kVect\), there exists a canonical map
  \begin{equation}\label{eq:canonical-coend-map}
    \int^{a \in \cat{C}}\int_{b \in \cat{C}} \cat{C}(b, a) \otimes T(a, b)
    \to \int_{b \in \cat{C}} \int^{a \in \cat{C}} \cat{C}(b, a) \otimes T(a, b).
  \end{equation}
  This follows from the fact that ends and coends are functorial;
  in particular, for any arrow \(f \from x \to y\), the following diagram commutes:
  \[
    \mathscale{0.92}{\begin{mytikzcd}[column sep=large, ampersand replacement=\&]
      {\int_b \cat{C}(b, x) \otimes T(y, b)} \&[10pt] {\int_b \cat{C}(b, y) \otimes T(y, b)} \\
      {\int_b \cat{C}(b, x) \otimes T(x, b)} \&[10pt] {\int^a\!\! \int_b \cat{C}(b, a) \otimes T(a, b)} \\
      \&\&[-30pt] {\int_b \int^a \cat{C}(b, a) \otimes T(a, b)}
      \arrow["{\int_b \cat{C}(b, f) \otimes T(y, b)}", from=1-1, to=1-2]
      \arrow["{\int_b \cat{C}(b, x) \otimes T(f, b)}"', from=1-1, to=2-1]
      \arrow["\mathrm{can}", from=1-2, to=2-2]
      \arrow["\mathrm{can}"', from=2-1, to=2-2]
      \arrow["{\exists!}", dotted, from=2-2, to=3-3]
      \arrow["{\int_b \mathrm{can}}"', curve={height=18pt}, from=2-1, to=3-3]
      \arrow["{\int_b \mathrm{can}}", curve={height=-18pt}, from=1-2, to=3-3]
    \end{mytikzcd}}
  \]
\end{remark}

\begin{lemma}[{\cite[p.~1]{day06:compac}}]\label{lem:day-yoneda}
  Let \(\cat{C}\) be a hom-finite category,
  \(P \from \cat{C}^{\op} \otimes_{\Bbbk} \cat{C} \to \kvect\) a functor,
  and suppose that the map given in \cref{eq:canonical-coend-map} is invertible.
  If for all \(a, b \in \cat{C}\) there is a natural isomorphism  \(\varphi \from \cat{C}(b, a) \iso \rd*{\cat{C}(a, b)}\),
  then
  \[
    \int^{a} P(a, a) \iso \int_{b} P(b, b).
  \]
\end{lemma}
\begin{proof}
  We calculate
  \begin{align*}
    \int^{a} P(a, a)
    & \cong \int^{a}\!\int_{b} \kvect(\cat{C}(a, b), P(a, b))
      \cong \int^{a}\!\int_{b}  \rd*{\cat{C}(a, b)}\otimes P(a, b) \\
    \overset{\varphi^{-1}}&{\cong} \int^{a}\!\int_{b}  \cat{C}(b, a)\otimes P(a, b)
                 \overset{\eqref{eq:canonical-coend-map}}{\cong}  \int_{b} \int^{a} \cat{C}(b, a) \otimes P(a, b)
                 \cong \int_{b} P(b, b).
  \end{align*}
  The first and last isomorphisms follow from \cref{lem:enriched-yoneda}
  and the second one is a consequence of \(\cat{C}(a, b)\) being finite\hyp{}dimensional and \(\kvect\) being rigid.
\end{proof}

In general, it seems difficult to decide whether a natural isomorphism as that in \cref{lem:day-yoneda} exists.
However, for certain classes of examples, trace maps provide us with viable candidates.

\begin{remark}\label{rem:day-morphism-invertible}
  Let \(\cat{C}\) be a hom-finite and pivotal category
  with pivotal structure \(\psi\from \rd{(\blank)}\to\ld{\!(\blank)}\),
  and suppose that \(\cat{C}(1,1) \cong \Bbbk\).
  For all \(a, b \in \cat{C}\), consider
  \begin{equation}\label{eq:day-canonical-candidate}
    \begin{aligned}
      \varphi_{a,b} \from \cat{C}(b, a) &\to \rd*{\cat{C}(a,b)} \\
      f &\mapsto \big(g \mapsto \tr(f \circ g) \defeq \ev^{\ell}_{a} \circ (\psi_{a}\otimes (f \circ g)) \circ\coev^{r}_{a}\big).
    \end{aligned}
  \end{equation}
  The dual of \(\varphi\) is given by
  \begin{align*}
    \rd*{(\varphi_{a,b})}\from \rrd*{\cat{C}(a,b)} \cong \cat{C}(a,b) &\to \rd*{\cat{C}(b,a)} \\
    g &\mapsto \big(f \mapsto \tr(f \circ g) = \tr(g \circ f) = \varphi_{b,a}(g)(f) \big).
  \end{align*}
  Thus, \(\varphi_{a,b}\) is injective if and only if \(\varphi_{b,a}\) is surjective.

  Suppose \(\Bbbk\) has characteristic zero and  \(\cat{C} \defeq \lmod{H}\) is the category of
  finite\hyp{}dimensional modules of a semisimple Hopf algebra \(H\).
  As a consequence of~\cite[Theorem~4]{larson-radford1988:SemisimpleCosemisimpleHopf},
  left and right duals coincide and we can choose the ``quantum trace'' of \cref{eq:day-canonical-candidate}
  to agree with the usual trace of endomorphisms between finite\hyp{}dimensional vector spaces.
  Let \(f\from M \to N\) be a morphism in \(\cat{C}\).
  By semisimplicity, the following short exact sequence splits:
  \[
    \begin{mytikzcd}[ampersand replacement=\&]
      0 \&\& {\ker f} \&\& M \&\& {\im f} \&\& 0
      \arrow[from=1-1, to=1-3]
      \arrow[from=1-3, to=1-5]
      \arrow["f", from=1-5, to=1-7]
      \arrow["\iota"{description}, curve={height=-18pt}, dashed, from=1-7, to=1-5]
      \arrow[, from=1-7, to=1-9]
    \end{mytikzcd}
  \]
  The arrow \(g = \iota \oplus 0 \from N \cong \im f \oplus N/\im f \to M\) satisfies
  \[
    \tr(fg) = \tr(f\iota)= \tr(\id_{\im f})=\dim \im f,
  \]
  implying that \(\varphi_{M,N}\) is injective.
\end{remark}

\begin{example}\label{ex:trace-map-not-invertible}
  Consider a field \(\Bbbk\) of characteristic \(p\), as well as \(q\in \mathbb{N}\) a multiple of \(p\).
  The group algebra \(H \defeq \Bbbk[\GL_q(p)]\) of the group of invertible \(q \times q\)\hyp{}matrices over \(\Bbbk\)
  provides us with an example of a pivotal hom-finite category where the morphism of \cref{eq:day-canonical-candidate} is not invertible.
  Matrix-vector multiplication turns \(\Bbbk^q\) into a simple \(H\)\hyp{}module whose endomorphism algebra is one-dimensional.
  This implies that for all \(f,g \in \End_{\GL_{q}(p)}(\Bbbk^p)\) there exists some \(\lambda \in \Bbbk\) such that \(\tr(fg)=\lambda \tr(\id_{\Bbbk^q}) = \lambda \cdot 0 = 0\).
\end{example}

\begin{remark}\label{rmk:reduction-of-assumptions}
  One can use a combination of \cref{cor:day-compact,lem:day-yoneda}
  to deduce that a given category of finite\hyp{}dimensional functors is right tensor representable.
  To reduce the number of axioms to be checked,
  note that the square of the left dualising functor being isomorphic to the identity is implied by the other assumptions.
  Let \((\cat{C}, d)\) be a hom-finite left \GV{} category with dualising functor \(D_{\ell}\),
  such that there are isomorphisms
  \(\cat{C}(a,b) \cong \rd*{\cat{C}(b, a)}\)
  and
  \(D_{\ell}(a \otimes b) \cong D_{\ell}b \otimes D_{\ell}a\),
  natural in \(a, b \in\cat{C}\).
  Using that
  \begin{align*}
    \cat{C}(d,a \otimes D_{\ell}b)
    \overset{\eqref{eq:left-gv-condition-and-inverse}}%
    &{\cong} \cat{C}(d \otimes D_{\ell}^{-1} (a \otimes D_{\ell} b), d)
      \cong \cat{C}(d \otimes b \otimes D_{\ell}^{-1}a, d) \\
    \overset{\eqref{eq:left-gv-condition-and-inverse}}%
    &{\cong} \cat{C}(D_{\ell}^{-1}a, D_{\ell}^{-1}(d \otimes b))
      \overset{\eqref{eq:gv:1:from:d:and:d:from:1}}%
      {\cong} \cat{C}(D_{\ell}^{-1}a, D_{\ell}^{-1}b)
      \cong \cat{C}(b,a),
  \end{align*}
  one obtains
  \[
    \cat{C}(a, D_{\ell}^2b)
    \overset{\eqref{eq:left-gv-condition-and-inverse}}%
    {\cong} \cat{C}(a \otimes D_{\ell}b, d)
    \!\cong\! \rd*{\cat{C}(d, a \otimes D_{\ell}b)}
    \!\cong\! \rd*{\cat{C}(b, a)}
    \!\cong\! \rrd*{\cat{C}(a,b)}
    \!\cong\! \cat{C}(a,b).
  \]
  The claim follows from the Yoneda lemma.
\end{remark}

\begin{lemma}\label{lem:semisimple-lets-ends-commute}
  Assume \(\Bbbk\) to be a perfect field.
  If the hom-finite category \(\cat{C}\) has finitely many objects
  and \(\Cpsh_{\fin}\)  is semisimple,
  then the canonical map
  \[
    \int^a\!\int_b T(b,b,a,a) \cong \int_b \int^a T(a,a,b,b)
  \]
  is invertible for all functors \(T \from \cat{C}^{\op}\otimes_{\Bbbk}\cat{C} \otimes_{\Bbbk}\cat{C}^{\op}\otimes_{\Bbbk}\cat{C} \to \kvect\).
\end{lemma}
\begin{proof}
  We endow the vector space \(A \defeq \bigoplus_{a, b\in\cat{C}}\cat{C}(a,b)\) with the structure of an associative unital algebra
  via the following multiplication, for \(f \in \cat{C}(a, b)\), \(g \in \cat{C}(c, d)\):
  \[
    f \cdot g \defeq
    \begin{cases*}
      g \circ f, & \text{if } b = c, \\
      0, & \text{otherwise.}
    \end{cases*}
  \]
  The unit \(\sum_{a \in\cat{C}} \id_a\) of \(A\)  is a sum of orthogonal idempotents.
  Thus, any right module \(M\) of \(A\) decomposes as a vector space into a direct sum \(M \cong \bigoplus_{a \in\cat{C}} M_a\), where \(M_a = M\cdot\id_a\), and the action of any  \(f \in\cat{C}(a,b)\)  defines a linear map \(M_a \to M_b\).
  Accordingly, a morphism of right \(A\)\hyp{}modules corresponds to a collection of homomorphisms \({\{\,\phi_a \from M_a \to N_a\,\}}_{a \in\cat{C}}\) such that
  \[
    \phi_b(m \cdot f) = \phi_a(m) \cdot f
    \qquad \qquad \text{for all \(m \in M_a\) and \(f\in \cat{C}(a,b)\)}.
  \]
  This defines a functor \(\Theta \from \rMod{A} \to \Cpsh\).
  Its quasi-inverse \(\Omega\)
  maps any \(F \in \Cpsh\) to the module \(\bigoplus_{a \in\cat{C}}Fa\),
  and any arrow \({\{\,\psi_a \from Fa \to Ga\,\}}_{a \in\cat{C}}\) to the module homomorphism \(\bigoplus_{a\in\cat{C}}\psi_a \from \bigoplus_{a \in\cat{C}} Fa \to \bigoplus_{a\in\cat{C}} Ga\).
  Thus, \(\Cpsh_{\fin}\) corresponds to the category \(\rmod{A}\) of finite\hyp{}dimensional right \(A\)\hyp{}modules.

  Since \(\Cpsh_{\fin}\) is semisimple, so is \(A\) and \(A^{\op}\).
  Furthermore, \(\Bbbk\) being perfect implies that the tensor product of any two finite\hyp{}dimensional semisimple algebras is semisimple, see~\cite[Section~3]{farb-dennis1993:Noncommutative}.
  In particular, we have that the algebra \(B \defeq A^{\op}\otimes_{\Bbbk}A \otimes_{\Bbbk}A^{\op}\otimes_{\Bbbk}A\) is semisimple,
  and
  \[
    [\cat{C}^{\op}\otimes_{\Bbbk}\cat{C} \otimes_{\Bbbk}\cat{C}^{\op}\otimes_{\Bbbk}\cat{C}, \kvect] \cong \rmod{B}.
  \]

  Write \(\cat{D} \defeq [\cat{C}^{\op}\otimes_{\Bbbk} \cat{C}, [\cat{C}^{\op} \otimes_{\Bbbk}\cat{C}, \kvect]]\);
  then the functor
  \[
    [\cat{C}^{\op}\otimes_{\Bbbk} \cat{C} \otimes_{\Bbbk}\cat{C}^{\op} \otimes_{\Bbbk}\cat{C}, \kvect] \to \cat{D},\qquad
    F \mapsto \widehat F,
  \]
  where \(\widehat{F}(x,y)[u,v] = F(u,v,x,y)\),
  is a \(\Bbbk\)\hyp{}linear equivalence of categories.

  Further, since limits commute with limits,
  the functor
  \[
    \mathrm{end}\from \cat{D} \to [\cat{C}^{\op} \otimes_{\Bbbk}\cat{C}, \kvect],
    \qquad
    F \mapsto (u,v \mapsto \Big(\int_b F(b,b)\Big)(u, v))
  \]
  is left exact.
  By semisimplicity, any short exact sequence in \(\cat{D}\) splits and split epimorphisms are preserved by all functors.
  Thus, \(\mathrm{end} \from\cat{D}\to [\cat{C}^{\op} \otimes_{\Bbbk}\cat{C}, \kvect]\) is exact and must preserve colimits.

  Given \(F\in  [\cat{C}^{\op}\otimes_{\Bbbk}\cat{C} \otimes_{\Bbbk}\cat{C}^{\op}\otimes_{\Bbbk}\cat{C}, \kvect]\), we now compute
  \begin{align*}
    \int^a\!\int_b F(a,a)(b,b)
    & \cong \int^a\!\int_b\widehat{F}(b,b)[a,a]
      \cong \int_b\int^a\widehat{F}(b,b)(a,a)\\
    & \cong \int_b\int^a\! F(a,a,b,b).\qedhere
  \end{align*}
\end{proof}

\subsection{Cauchy completions}\label{sec: trep-rigid-fgp}

\textsc{Having assembled all of the necessary tools},
we can now state explicit criteria for the Cauchy completion of (the opposite of) a \(\Bbbk\)\hyp{}linear category to carry certain duality structures.
As with modules over (commutative) rings, these notions are closely connected with objects being finitely\hyp{}generated projective.
Again, we implicitly assume all categories and functors to be \(\Bbbk\)\hyp{}linear, for a field \(\Bbbk\).
However, we do not require finiteness of hom-spaces and work solely under the assumptions made in \cref{hypo:eso-enriched}.

\begin{definition}\label{def:k-linear-cauchy-completion}
  Suppose that \(\cat{C}\) is a \(\Bbbk\)\hyp{}linear category.
  The \emph{Cauchy completion} of \(\cat{C}\) is the full subcategory \(\overline{\cat{C}}\) of \(\psh\)
  that consists of all presheaves \(F\),
  such that \(\psh(F, \blank)\) commutes with all small limits.
\end{definition}

\begin{remark}
  \Cref{def:k-linear-cauchy-completion} differs from the usual definition of the Cauchy completion of a \(\Bbbk\)\hyp{}linear category \(\cat{C}\)
  as the additive and idempotent completion of \(\cat{C}\).
  However, by~\cite[Proposition~2]{borceux-dejean1986:CauchyCompletionCategoryTheory}
  and~\cite[Corollary~4.22]{lack-tendas2022:FlatFilteredColimts},
  these two notions are equivalent.
  We choose the former because it is more convenient to work with for the purposes of this thesis.
\end{remark}

\begin{example}\label{ex:cauchy-tensor-prod-in-prof-fin}
  Let \(\iota\from \cat{C} \emb \overline{\cat{C}}\)
  denote the inclusion of a category \(\cat{C}\) into its Cauchy completion,
  sending \(x \in \cat{C}\) to \(\cat{C}(\blank, x)\).
  There is an adjoint equivalence
  \[
    \big(\overline{\cat{C} }(\blank, \iota {   =}),\, \overline{\cat{C} }(\iota \blank, { = }),\, \eta,\, \varepsilon\big),
  \]
  between \(\cat{C}\) and its Cauchy completion in the monoidal bicategory of profunctors,
  see \cref{MonoidalBicatProf}.
  The unit
  \[
    \eta_{x, y} \from \cat{C}(x, y) \to \int^{\overline{c} \in \overline{\cat{C}}} \overline{\cat{C}}(\iota_x, \overline{c}) \times \overline{\cat{C} }(\overline{c}, \iota_y)
  \]
  is defined by
  \[
    \cat{C}(x, y)
    \xrightarrow{\;\iota_{x, y}\;} \overline{\cat{C} }(\iota_{x}, \iota _{y})
    \cong \int^{\overline{c} \in \cat{C} } \overline{\cat{C} }(\iota_x, \overline{c}) \times \overline{\cat{C} }(\overline{c}, \iota_y),
  \]
  and for the counit we have
  \begin{align*}
    \varepsilon_{\overline{x}, \overline{y}}
    \from \int^{c \in \cat{C} } \overline{\cat{C}}(\overline{x}, \iota_c) \times \overline{\cat{C} }(\iota_c, \overline{y})
    &\to \overline{\cat{C} }(\overline{x}, \overline{y}) \\
    [(f, g)] &\mapsto g \circ f.
  \end{align*}

  Since \(\iota\) is fully faithful,
  it immediately follows that \(\eta_{x,y}\) is an isomorphism, for all \(x, y \in \cat{C}\),
  and the counit \(\varepsilon\) is one by~\cite[Theorem~1.1]{adamek01}.
\end{example}

In principle, the absence of free objects in the category of copresheaves necessitates a development of these notions in abstract categorical terms;
this is discussed extensively for example in~\cite{prest09:purit}.
For our purposes, the following characterisation in terms of the Cauchy completion of \(\cat{C}^{\op}\) suffices.

\begin{lemma}\label{lem:projective-if-representable}
  Let \(\cat{C}\) be a category.
  For any \(F \in \Cpsh\), the following are equivalent:
  \begin{thmlist}[noitemsep]
  \item\label{itm:fgp-char-fgp} \(F\) is finitely\hyp{}generated projective,
  \item\label{itm:fgp-char-dir-sum} \(F\) is a direct summand of a finite direct sum of representable functors, and
  \item\label{itm:fgp-char-colim} the functor \(\Cpsh(F, \blank)\) commutes with small colimits.
  \end{thmlist}
\end{lemma}
\begin{proof}
  The equivalence between~\ref{itm:fgp-char-fgp} and~\ref{itm:fgp-char-dir-sum} is proven in Corollary~10.1.14 of~\cite{prest09:purit}.
  In order to show that~\ref{itm:fgp-char-dir-sum} and~\ref{itm:fgp-char-colim} are equivalent,
  observe that the full subcategory \(\overline{\cat{C}^{\op}}\) of \(\Cpsh\) consisting of direct summands of finite direct sums of representable functors is Cauchy complete by~\cite[Corollary~4.22]{lack-tendas2022:FlatFilteredColimts}.
  The claim now follows by proceeding analogous to Proposition~2 of~\cite{borceux-dejean1986:CauchyCompletionCategoryTheory}.
\end{proof}

\begin{notation}
  In view of \cref{lem:projective-if-representable}, we shall adopt the following notation
  to emphasise that \(X \in \overline{\cat{C}^{\op}}\) is a direct summand of a direct sum of representables:
  \[
    \begin{mytikzcd}[cramped, ampersand replacement=\&]
      X \& {\bigoplus\limits_{i=1}^n \yo_{u_i},}
      \arrow["{\iota_X}", shift left, from=1-1, to=1-2]
      \arrow["{\pi_X}", shift left, from=1-2, to=1-1]
    \end{mytikzcd}
  \]
  where \(\pi_X \circ \iota_X = \id_X\), and
  \[
    \yo\from \cat{C}^{\op} \to [{(\cat{C}^{\op})}^{\op}, \kVect] = [\cat{C}, \kVect], \qquad x \mapsto \cat{C}^{\op}(\blank, x) = \cat{C}(x, \blank)
  \]
  is the contravariant Yoneda embedding as in \cref{eq:yoneda-embedding}.
\end{notation}

\begin{proposition}\label{prop:cauchy-completion-is-trep}
  For all objects \(
  \begin{mytikzcd}[cramped, ampersand replacement=\&]
    X \& {\bigoplus_{i=1}^n \yo_{u_i}}
    \arrow["{\iota_X}", shift left, from=1-1, to=1-2]
    \arrow["{\pi_X}", shift left, from=1-2, to=1-1]
  \end{mytikzcd}
  \) and
  \(
  \begin{mytikzcd}[cramped, ampersand replacement=\&]
    Y \& {\bigoplus_{j=1}^m \yo_{v_j}}
    \arrow["{\iota_Y}", shift left, from=1-1, to=1-2]
    \arrow["{\pi_Y}", shift left, from=1-2, to=1-1]
  \end{mytikzcd}
  \),
  the right internal hom
  of the Cauchy completion \(\overline{\cat{C}^{\op}}\)
  of a right closed monoidal category \(\cat{C}^{\op}\)
  exists and satisfies
  \begin{equation}\label{eq:char-of-internal hom-of-trep}
    \begin{mytikzcd}[ampersand replacement=\&]
      {[X,Y]}_r \& \& {[\oplus_{i=1}^n \yo_{u_i}, \oplus_{j=1}^m \yo_{v_j}]\cong \bigoplus\limits_{i=1}^n \bigoplus\limits_{j=1}^m \yo_{{[u_i,v_j]}_r}.}
      \arrow["{{[\pi_{X},\iota_Y]}_r}", shift left, from=1-1, to=1-3]
      \arrow["{{[\iota_X, \pi_Y]}_r}", shift left, from=1-3, to=1-1]
  \end{mytikzcd}
  \end{equation}
\end{proposition}
\begin{proof}
  The category \(\overline{\cat{C}^{\op}} \subseteq \Cpsh\) is closed under taking tensor products
  since the Yoneda embedding is strong monoidal
  and \(\overline{\cat{C}^{\op}}\) is the full subcategory of \(\Cpsh\) consisting of direct summands of finite direct sums of representables.

  Let \(X, Y \in \overline{\cat{C}^{\op}}\).
  Write \(X = \colim_i\cat{C}(u_i, \blank)\) and \(Y = \colim_j\cat{C}(v_j, \blank)\) as direct summands of finite direct sums of representables.
  We compute
  \begin{align*}
    {[X,Y]}_r&
              \overset{\eqref{eq:Day-internal-hom-right}}{=} \int_{a, b}\!\kVect(\cat{C}(a \otimes \blank, b), \kVect(Xa, Yb))
              \overset{\eqref{eq:coYoneda-covariant}}{\cong} \int_{a}\!\!\kVect(Xa, Y(a\otimes \blank))\\
            &\hspace{-0.5cm} \cong \int_{a} \kVect(\colim_i\cat{C}(u_i, a), \colim_j \cat{C}(v_j, a \otimes \blank)) \\
            &\hspace{-0.5cm} \cong \lim\nolimits_i \colim_j  \int_{a}\kVect(\cat{C}(u_i, a), \cat{C}(v_j, a\otimes \blank))
              \cong \lim\nolimits_i \colim_j\cat{C}(v_j, u_i \otimes \blank) \\
            &\hspace{-0.5cm} \cong \lim\nolimits_i \colim_j\cat{C}({[u_i, v_j]}_r, \blank).
  \end{align*}
  It follows that \(\overline{\cat{C}^{\op}}\) is right closed monoidal and that \cref{eq:char-of-internal hom-of-trep} holds.
\end{proof}

\begin{example}
  \Cref{prop:cauchy-completion-is-trep} can be understood as a variation of the fact that
  a direct summand or direct sum of (rigidly) dualisable objects is dualisable.
  Indeed, suppose \(\cat{C}^{\op}\) to be a closed monoidal category
  and consider three direct summands
  \(\begin{mytikzcd}[ampersand replacement=\&, cramped]
    X \& {U}
    \arrow["{\iota_X}", shift left, from=1-1, to=1-2]
    \arrow["{\pi_X}", shift left, from=1-2, to=1-1]
  \end{mytikzcd}
  \),
  \(
  \begin{mytikzcd}[ampersand replacement=\&, cramped]
    Y \& V
    \arrow["{\iota_Y}", shift left, from=1-1, to=1-2]
    \arrow["{\pi_Y}", shift left, from=1-2, to=1-1]
  \end{mytikzcd}
  \), and
  \(\begin{mytikzcd}[ampersand replacement=\&, cramped]
    Z \& {W}
    \arrow["{\iota_W}", shift left, from=1-1, to=1-2]
    \arrow["{\pi_W}", shift left, from=1-2, to=1-1]
  \end{mytikzcd}\) of objects in \(\overline{\cat{C}^{\op}}\).
  The following diagram,
  whose horizontal arrows are the isomorphisms of the tensor--hom adjunction of \(\overline{\cat{C}^{\op}}\),
  commutes:
  \[
    \begin{mytikzcd}[ampersand replacement=\&]
      {\overline{\cat{C}^{\op}}(X * Y, Z)} \&\&\& {\overline{\cat{C}^{\op}}(Y,{[X,Z]}_{r})} \\
      \\
      {\overline{\cat{C}^{\op}}(U * V, W)} \&\&\& {\overline{\cat{C}^{\op}}(V, {[U,W]}_{r})}
      \arrow["{{\phi_{X,Y,Z}}}", from=1-1, to=1-4]
      \arrow["{{\overline{\cat{C}^{\op}}(\pi_X * \pi_Y, \iota_W)}}"', shift right=2, from=1-1, to=3-1]
      \arrow["{{\overline{\cat{C}^{\op}}(\pi_V,{[\pi_X,\iota_W]}_r)}}"', shift right=2, from=1-4, to=3-4]
      \arrow["{{\overline{\cat{C}^{\op}}(\iota_X * \iota_Y, \pi_Z)}}"', shift right=2, from=3-1, to=1-1]
      \arrow["{{\phi_{U,V,W}}}", from=3-1, to=3-4]
      \arrow["{{\overline{\cat{C}^{\op}}(\iota_V,{[\iota_X,\pi_Z]}_r)}}"', shift right=2, from=3-4, to=1-4]
    \end{mytikzcd}
  \]

  Thus the unit and counit of the adjunction
  \(\adj{X * \blank}{{[X, \blank]}_r}{\overline{\cat{C}}}{\overline{\cat{C}}}\)
  satisfy
  \begin{equation}\label{eq:unit-counit-formula-direct-summand}
    \begin{aligned}
      \eta^{(X)}_Y
      &= Y \xrightarrow{\ \iota_Y\ } V
      \xrightarrow{\,\eta_V^{(U)}\,} [U, U * V]
        \xrightarrow{[\iota_X, \pi_X * \pi_Y]} [X, X * V], \\
      \varepsilon^{(X)}_Y
      &= X * [X, Y]
        \xrightarrow{\iota_X* [\pi_X,\iota_Y]} U * [U, V]
        \xrightarrow{\,\varepsilon^{(U)}_V\,} V
        \xrightarrow{\ \pi_Y\ } Y.
    \end{aligned}
  \end{equation}
\end{example}

\begin{corollary}\label{cor:pres-ref-dual-struct}
  Let \(\cat{C}^{\op}\) be a right closed monoidal category.
  We have:
  \begin{thmlist}[noitemsep]
    \item \(\cat{C}^{\op}\) is right rigid if and only if \(\overline{\cat{C}^{\op}}\) is.\label{itm:rigid-implies-rigid}
    \item \(\cat{C}^{\op}\) is right tensor representable if and only if \(\overline{\cat{C}^{\op}}\) is.\label{itm:trep-implies-trep}
    \item \((\cat{C}^{\op}, d)\) is a right \GV{} category if and only if \((\overline{\cat{C}^{\op}}, \yo_{d})\) is.\label{itm:GV-implies-GV}
  \end{thmlist}
\end{corollary}
\begin{proof}
  That right rigidity and tensor representability of \(\cat{C}^{\op}\) imply the same property for \(\overline{\cat{C}^{\op}}\)
  follows from the description of the internal hom of \(\overline{\cat{C}^{\op}}\)
  and the units and counits of the tensor-hom adjunctions,
  see \cref{eq:unit-counit-formula-direct-summand}.

  If \((\cat{C}^{\op}, d)\) is a right \GV{} category,
  then its right dualising functor is, up to natural isomorphism, given by \({[\blank, d]}_r \from\cat{C} \to\cat{C}^{\op}\).
  A direct computation using \cref{prop:cauchy-completion-is-trep} shows that
  \({[\blank, \yo_{d}]}_r \from \overline{\cat{C}^{\op}}^{\op}\to \overline{\cat{C}^{\op}}\) is an equivalence,
  and therefore \(\big(\overline{\cat{C}^{\op}}, \yo_{d}\big)\) is right \GV{}.

  By \cref{prop:cauchy-completion-is-trep} the right internal hom of two representable functors \(\yo_{x}\) and \(\yo_{y}\) is \({[\yo_{x},\yo_{y}]}_r \cong\yo_{{[x,y]}_r}\).
  Thus, the converse of any of the three statements is a consequence of the fact that, via the Yoneda embedding, \(\cat{C}^{\op}\) is equivalent as a right closed monoidal category to the full subcategory of \(\overline{\cat{C}^{\op}}\) whose objects are representable functors.
\end{proof}

\begin{remark}\label{rmk:fgp-is-trep-but-not-rigid}
  The linearisation \(\Bbbk\cat{W}\) of the tensor representable category \(\cat{W}\) discussed in \cref{thm:props-of-c} is tensor representable,
  with left and right dualising functors \(L\) and \(R\),
  but not rigid.
  The latter follows from the fact that by the proof of \cref{thm:props-of-c} there exists an object \(t \in\cat{W}\)
  whose canonical morphism \({[t,1]}_r\otimes t \to {[t,t]}_r\) is not invertible.

  Since the right dualising functor \(R\) of \(\cat{W}\) is an anti\hyp{}equivalence,
  and so \(L \cong R^{-1}\),
  there also is a tensor representable, but not rigid, structure on \(\Bbbk\cat{W}^{\op}\).
  One takes the left dualising functor to be \(R\), and the right one to be \(L\).
  By \cref{cor:pres-ref-dual-struct}, \(\overline{\Bbbk\cat{W}^{\op}}\)
  is tensor representable but not rigid.
\end{remark}

Notice that this is in stark contrast to many cases arising in representation theory,
like modules over commutative rings or \(\Bbbk\)\hyp{}algebras,
where rigidity and tensor representability are equivalent;
see for example~\cite[Proposition~2.1]{niefield17:coexp} for a slightly more general statement.
Since a \(\Bbbk\)\hyp{}algebra can be interpreted as a \(\Bbbk\)\hyp{}linear category with one object,
the following proposition may be seen as a ``many object'' version of the classical case.

\begin{proposition}\label{prop:characterisation-of-rigidity}
  Let \(\cat{C}\) be a left rigid monoidal category.
  The following are equivalent for some \(X \in \Cpsh\):
  \begin{thmlist}[noitemsep]
  \item\label{itm:char-rigid} it has a right rigid dual,
  \item\label{itm:char-trep} there exists a \(\mathbb{D}X \in \Cpsh\) such that \(\adj{X * \blank}{\mathbb{D}X * \blank}{\Cpsh}{\Cpsh}\), and
  \item\label{itm:char-fgp} it is finitely\hyp{}generated projective.
  \end{thmlist}
\end{proposition}
\begin{proof}
  By definition, \(\text{\ref{itm:char-rigid}}\!\implies\!\text{\ref{itm:char-trep}}\).
  To show that \(\text{\ref{itm:char-trep}}\!\implies\!\text{\ref{itm:char-fgp}}\),
  we assume that there exists a \(\mathbb{D}X \in \Cpsh\)
  such that  \(\adj{X * \blank}{\mathbb{D}X * \blank}{\Cpsh}{\Cpsh}\),
  and fix a small colimit \(\colim_{i \in I} Fi \in \Cpsh\) of some diagram \(F \from I \to  \Cpsh\).
  For every \(i \in I\), write \(\iota_i \from Fi \to \colim_{i \in I} Fi \in \Cpsh\) for its structure morphisms.
  Now consider the commutative diagram:
  \[
    \begin{mytikzcd}[ampersand replacement=\&,column sep=3.5em]
      {\colim_{i \in I} \Cpsh (X, Fi)} \&\& {\Cpsh (X, \colim_{i \in I} Fi)} \\
      \\
      {\colim_{i \in I} \Cpsh (1,\mathbb{D}X * Fi)} \&\& { \Cpsh (1, \mathbb{D}X * \colim_{i \in I} Fi)}
      \arrow["{\colim_{i \in I}(\Cpsh(X,\iota_i))}", from=1-1, to=1-3]
      \arrow["{\colim_{i \in I} \phi_{X, Fi}}", from=1-1, to=3-1]
      \arrow["{\phi_{X, \colim_{i \in I} Fi}}", from=1-3, to=3-3]
      \arrow["{\colim_{i \in I}(\Cpsh(1,\mathbb{D}X * \iota_i))}", from=3-1, to=3-3]
    \end{mytikzcd}
  \]
  Its horizontal arrows are induced by the universal property of the colimit and the vertical arrows are due to the tensor--hom adjunction of \(\Cpsh\).
  The functors \(\Cpsh(1, \blank)\) and \(\mathbb{D}X \otimes \blank\) commute with all small colimits;
  the first one due to the fact that \(\Cpsh(1, \blank)\) is finitely\hyp{}generated projective,
  see \cref{lem:projective-if-representable}, and the second one due to Day convolution being a colimit.
  Therefore, the horizontal arrow at the bottom is invertible.
  Since the vertical arrows are also invertible,
  the canonical arrow \(\colim_{i \in I} \Cpsh (X, Fi) \to \Cpsh (X, \colim_{i \in I} Fi)\)
  displayed at the top of the diagram must be an isomorphism.
  Again, using \cref{lem:projective-if-representable},
  \(X\) must be finitely\hyp{}generated projective.
  Thus,~\ref{itm:char-trep} implies~\ref{itm:char-fgp}.

  Finally, if \(X \in \Cpsh\) is finitely\hyp{}generated projective, then it is contained in the Cauchy-completion of \(\cat{C}^{\op}\),
  which is right rigid.
  Then, by \cref{cor:pres-ref-dual-struct}, so is \(\overline{\cat{C}^{\op}}\).
  Therefore \(X\) admits a rigid dual and \(\text{\ref{itm:char-fgp}}\!\implies\!\text{\ref{itm:char-rigid}}\).
\end{proof}

\section{Applications}\label{sec:mky-2-grps}

\textsc{We conclude the chapter by discussing} several examples.

We investigate Boolean algebras,
their applications in group and ring theory,
and how they induce abelian \(\Bbbk\)\hyp{}linear \GV{} categories.

Next, we focus on Mackey functors.
These are, roughly speaking, collections of vector spaces indexed by all subgroups of a fixed finite group,
together with morphisms subject to relations resembling the behaviour of
induction, restriction, and conjugation operations, including the eponymous Mackey identity;
see~\cite{lindner76:mackey,thevenaz95:mackey}.
Finite\hyp{}dimensional Mackey functors form a \GV{} category~\cite{panchadcharam07:mackey};
we show in \cref{prop:mky-GV-and-rigid} that it is rigid if and only if it is semisimple.

The last example arises in the study of crossed modules,
which—in categorical terms—correspond to strict 2-groups.
The functors from any finite strict 2-group to \(\kvect\) form an abelian monoidal category
that is equivalent to a direct sum of representation categories of the isotropy-group of the monoidal unit of the 2-group.
In \cref{prop:2-grp-GV-and-rigid} we prove the rigidity of this category to be equivalent to the semisimplicity of a certain group algebra.

\subsection{Boolean algebras}\label{sec:boolean-algebras}

\begin{definition}\label{def:lattice}
  A \emph{lattice} consists of a set \(L\) together with two associative and commutative operations
  \[
    \land \from L\times L \to L, \quad (a,b) \mapsto a \land b
    \quad\text{and}\quad
    \lor \from L \times L \to L, \quad (a,b) \mapsto a \lor b,
  \]
  called \emph{meet} and \emph{join}, which satisfy the \emph{absorption laws}, for all \(a, b \in L\):
  \[
    a \lor ( a\land b) = a,\qquad\qquad
    a \land (a \lor b) = a.
  \]
\end{definition}

\begin{remark}
  Any lattice \((L, \land, \lor)\) defines a poset via the relation
  \[
    a \leq b \iff b = a \lor b \iff a = a \land b, \qquad\qquad a, b \in L.
  \]

  Conversely, a poset \((P,\leq)\) that admits for any pair of objects \(a,b \in P\) a least upper bound \(a\lor b \in P\) and a greatest lower bound \(a \land b\in P\) is a lattice.
\end{remark}

A direct computation shows that an element \(1 \in L\) is maximal with respect to the partial order of the lattice \(L\) if and only if \(1 \lor a = 1\)  for all \(a \in L\).
Analogously, \(0 \in L\) being minimal equates to \(0\land a = 0\) for any \(a\in L\).
Minimal and maximal elements are unique.
In case they exist, we call \(L\) \emph{bounded}.

\begin{definition}\label{def:Boolean-algebra}
  A \emph{Boolean algebra} is a bounded lattice \((L, \land, \lor)\) satisfying the following \emph{distributivity}
  condition for all \(a,b,c\in L\):
  \begin{equation}\label{eq:distributivity-equation}
    a \land (b \lor c) = (a \land b) \lor (a \land c),\qquad
    a \lor (b \land c) = (a \lor b) \land (a \lor c),
  \end{equation}
  and every \(a\in L\) admits a \emph{complement} \(a^{\perp}\in L\) in the sense that
  \[
    a \lor a^{\perp} = 1
    \qquad\text{and}\qquad
    a \land a^{\perp} = 0.
  \]
\end{definition}

\begin{remark}
  For any Boolean algebra \((L, \land, \lor)\),
  any of the two distributivity requirements of \cref{eq:distributivity-equation} implies the other.
  Further, a direct computation shows that complements are unique.
  We obtain an involutive map \({( \blank )}^{\perp} \from L \to L\) that maps \(a\) to \(a^{\bot}\).
  The maximal element \(1 \in L\) is a unit for \(\land\):
  \[
    a \land 1 = a \land (a \lor a^{\bot}) = a, \qquad \text{for all } a \in L.
  \]

  An analogous computation shows that the minimal element \(0\) in a Boolean algebra satisfies \(a \lor 0 = a\).
\end{remark}

\begin{example}\hfill
  \begin{itemize}
    \item\emph{Central idempotents}:
    the set \(C \defeq \{\,e\in Z(R)\mid e^2=e\,\}\) of central idempotents of a ring \(R\) is a Boolean algebra
    when endowed with
    \begin{align*}
      &\land \from C\times C \to C &&e \land f = ef,\\
      &\lor \from C \times C \to C &&e \lor f = e+f-ef,\\
      &{( \blank )}^{\perp}\from C \to C &&e^{\perp}=1-e.
    \end{align*}
    \item\emph{Annihilators in semiprime rings}:
    consider a commutative semiprime ring \(R\) with \(1\).
    That is, its Jacobson radical is trivial.
    In case \(R\) is furthermore Artinian, this is equivalent to it being semisimple.
    The \emph{annihilator} of an ideal \(I \subset R\) is \(I^{\perp} \defeq \{\,x\in R \mid xI=0\,\}\);
    it is a radical ideal.
    We define on the set \(\mathrm{Ann}(R)\) of annihilators the maps
    \begin{align*}
      \land \from \mathrm{Ann}(R) \times \mathrm{Ann}(R)\to \mathrm{Ann}(R), \qquad I \land J = I \cap J, \\
      \lor \from \mathrm{Ann}(R) \times \mathrm{Ann}(R)\to \mathrm{Ann}(R), \qquad I \lor J = {(I + J)}^{\perp}.
    \end{align*}
    This defines the structure of a Boolean algebra on \(\mathrm{Ann}(R)\).
    In fact,
    the (right) annihilators of a not necessarily commutative ring
    form a Boolean algebra
    if and only if the ring is semiprime, see~\cite{dube-taherifar2021:LatticeAnnihilatorIdeasSemiprime}.
    \item \emph{The subgroup lattice}:
    let \(H \subseteq G\) be a subgroup of a finite group.
    We write \(H^{\perp}\) for the intersection of all maximal subgroups that do not contain \(H\).
    The minimal group constructed in this manner is the Frattini subgroup \(\Phi(G)=G^{\perp}\) of \(G\).
    Deaconescu, Isaacs, and Wall showed that the set
    \[
      \{\, \Phi(G) \subseteq H \subseteq G \mid G= HH^{\perp} \,\},
    \]
    partially ordered under inclusion forms a Boolean algebra,
    see~\cite{deaconescu-isaacs-walls2011:BooleanLatticeFrattiniSubgroup}.
  \end{itemize}
\end{example}

Barr showed in~\cite{barr79} that the duality of Boolean algebras fits within the framework of \GV{} categories.

\begin{proposition}\label{prop:complement-lattice-form-GV-category}
  Let \((L, \land, \lor)\) be a Boolean algebra.
  The associated poset-category \(\cat{L}\) is a left \GV{} category
  with meet as tensor product, \(1\) as monoidal unit, and \(0\) as dualising object.
  The dualising functor is induced by \({( \blank )}^{\perp}\).

  If \(L\) is finite,
  the category \(\langle\cat{L}, \kvect\rangle\)
  of ordinary functors from \(\cat{L}\)
  to the category of finite\hyp{}dimensional \(\Bbbk\)\hyp{}vector spaces
  can be equipped with a right \GV{} structure,
  with \(\rd*{\Bbbk\cat{L}( \blank , s)}\) as dualising object and dualising functor
  \[
    \mathfrak{D}_r \from \langle \cat{L}, \kvect\rangle^{\op} \to \langle \cat{L}, \kvect\rangle, \qquad
    F \mapsto {F(\blank^{\perp})}^{*}.
  \]
\end{proposition}
\begin{proof}
  A direct consequence of \cref{eq:distributivity-equation} is that for \(a \leq c\) and \(b \leq d\), we have
  \(a \land b \leq c \land d\).
  Thus, since \(\land \from L \times L \to L\) is associative, commutative, and unital, it induces a (symmetric) monoidal structure on \(\cat{L}\).

  To see that \((\cat{L}, 0)\) is a  \GV{} category,
  we compute its left internal hom.
  Let us fix a triple of elements \(a,b,c \in L\).
  We have
  \begin{align*}
    b \land a \leq c
    & \iff (b \land a) \lor a^{\bot} \leq c \lor a^{\bot}
      \iff (b \lor a^{\bot}) \land (a \lor a^{\bot}) \leq c \lor a^{\bot} \\
    & \iff b \lor a^{\bot} \leq c \lor a^{\bot}
      \iff b \leq c \lor a^{\perp},
  \end{align*}
  where the last equivalence is due to the fact that \(b = b \lor 0 \leq b \lor a^{\bot}\)
  and \(b \leq c \lor a^{\perp}\), implying that \(b\lor a^{\perp} \leq (c \lor a^{\perp})\lor a^{\perp}= c \lor a ^{\perp}\).
  It follows that
  \[
    {[\blank, \bblank]}_{\ell} \from L \times L \to L, \qquad {[a, b]}_{\ell} \defeq {(a \land b^{\bot})}^{\bot} = b \lor a^{\perp},
  \]
  defines the left internal hom of \(\cat{L}\),
  and the order reversing involution
  \[
    {[\blank, 0]}_{\ell} = {( \blank )}^{\perp} \from L \to L
  \]
  induces an equivalence of categories \(D_{\ell} \from \cat{L}^{\op}\to\cat{L}\).
  As discussed in \cref{rmk:gv-d-from-D},
  \((\cat{L}, 0)\) is a left \GV{} category.

  By \cref{rmk:linearisation}, we have an equivalence \(\langle \cat{L}, \kvect\rangle \cong [\Bbbk\cat{L}, \kvect]\).
  Under the assumption that \(L\) is finite,
  \cref{prop:fun-cat-GV} shows that \(\langle \cat{L}, \kvect\rangle\) has the structure of a right \GV{} category
  with the specified dualising object and dualising functor.
\end{proof}

Suppose \(L\) is a finite Boolean algebra and \(\cat{L}\) is its poset-category.
As discussed in the first steps of the proof of \cref{lem:semisimple-lets-ends-commute},
the category \([\Bbbk\cat{L}, \kvect]\) can be identified with the finite\hyp{}dimensional right modules
of the path algebra \(A \cong \bigoplus_{e,f \in L}\Bbbk\cat{L}(e,f)\) of the poset associated to \(L\).

\begin{example}\label{ex:upper-triangular-subalgebra}
  The partial order on the set \(L \defeq \{\,0,a,b,1\,\}\) displayed in the following diagram defines a Boolean algebra:
  \[
    \begin{mytikzcd}[ampersand replacement=\&]
      \& a \\
      0 \&\& 1 \\
      \& b
      \arrow["v", from=1-2, to=2-3]
      \arrow["u", from=2-1, to=1-2]
      \arrow["x"', from=2-1, to=3-2]
      \arrow["y"', from=3-2, to=2-3]
    \end{mytikzcd}
  \]
  A direct computation shows that its path algebra \(A\) is isomorphic to
  a subalgebra of the \(\Bbbk\)\hyp{}valued \(4 \times 4\) upper triangular matrices:
  \[
    A \cong\left\{
      \begin{pmatrix}
        0 & u & x & z \\
        0 & a & 0 & v \\
        0 & 0 & b & y \\
        0 & 0 & 0 & 1
      \end{pmatrix} \;\middle| \;
      0,u,x,z,a,v,b,y,1 \in \Bbbk
    \right\}.
  \]
\end{example}

Subalgebras of upper triangular matrices were studied by Thrall in~\cite{thrall1948:GeneralisationsQuasiFrobenius}
to provide generalisations of quasi\hyp{}Frobenius algebras.\note{%
  A \(\Bbbk\)\hyp{}algebra is \emph{quasi\hyp{}Fro\-be\-ni\-us} if all projective right modules are injective.%
}
The following definition generalises this property.

\begin{definition}\label{def:qf-2}
  A \(\Bbbk\)\hyp{}algebra is called \textsc{qf}-2 if
  all indecomposable projective right and all
  indecomposable projective left modules
  have simple socles.
\end{definition}

Using the \GV{} structure of \cref{prop:complement-lattice-form-GV-category},
we will show that the path algebra of any finite Boolean algebra \(L\) is \textsc{qf}-2.
Hereto we need the following observation.
Taking complements in \(L\) extends to an anti\hyp{}algebra isomorphism \(\phi \from A \to A\);
its pushforward \({( \blank )}_{\phi} \from \rmod{A} \to \lmod{A}\) is an involutive equivalence of categories.
It maps any left module \(M\) to the right module \(M_{\phi}\)
that has the same underlying vector space
and is endowed with the action \(a \lact m  \defeq m \ract \phi(a)\) for all \(m\in M\) and \(a\in A\).

\begin{proposition}\label{prop:qf2-algebras-from-boolean}
  Let \(A\) be the path algebra of a finite Boolean algebra \(L\).
  Then
  \[
    \rd*{( \blank_{\phi} )} \from {(\rmod{A})}^{\op} \to \rmod{A}
  \]
  is an equivalence of categories.
  In particular, \(M \in \rmod{A}\) is projective if and only if \(\rd*{M_{\phi}}\) is injective.
  Further, \(A\) is \textsc{qf}-2, and quasi\hyp{}Frobenius if and only if \(|L| = 1\).
\end{proposition}
\begin{proof}
  The first statement follows directly from the equivalence of categories between \([\Bbbk\cat{L}, \kvect]\) and \(\rmod{A}\)
  given in the proof of \cref{lem:semisimple-lets-ends-commute},
  and the definition of the dualising functor of \([\Bbbk\cat{L}, \kvect]\) stated in
  \cref{prop:complement-lattice-form-GV-category}.

  In order to show the second claim,
  we observe that the finite\hyp{}dimensional algebra \(A = \bigoplus_{n \geq 0} A_n\) is graded by the path lengths.
  The elements of \(L\) form a basis of \(A_0\) and correspond to the \emph{primitive idempotents}%
  —non-zero idempotents, which cannot be written as a sum of two non-zero orthogonal idempotents.
  Furthermore, the Jacobson radical of \(A\) is \(J(A) = \bigoplus_{n\geq 1} A_n\).
  The indecomposable projectives of \(A\) are of the form \(eA\), for \(e\in L\).
  Note that \(eA\) has a vector space basis by paths \([e,y]\) for \(e \leq y \leq 1\).
  The socle of \(eA\) is
  \[
    \mathrm{soc}(eA)
    = \{\,m \in eA \mid mJ(A)=0\,\}
    = \mathrm{span}_{\Bbbk}\{[e,1]\},
  \]
  which is one-dimensional and therefore simple.
  Using that \(\lmod{A} \cong \rmod{A}\), it follows that \(A\) is a \textsc{qf}-2 algebra.

  For \(|L|=1\), we obtain the trivial (quasi-)Frobenius-algebra \(A=\Bbbk\).
  Otherwise, there exists an \(e\in L\) such that \(\dim eA\geq 2\).
  A direct computation shows that \(\mathrm{soc}(eA)\cong 1 A\).
  If \(eA\) was injective,
  the inclusion \(1 A \cong \mathrm{soc}(eA) \longhookrightarrow eA\) would have a retraction,
  contradicting the indecomposability of \(eA\).
\end{proof}

\subsection{Mackey functors}\label{sec:mackey-functors}

\textsc{Let \(G\) be a finite group} and write \(\Spiso{G}\) for the category of isomorphism classes of spans of finite \(G\)\hyp{}sets.%
\marginnote{%
  In this section we again implicitly assume all categories and functors to be \(\Bbbk\)\hyp{}linear.
}
One can show that each hom-set is a free and finitely\hyp{}generated commutative monoid,
see for example the discussion preceding Lemma~2.1 of~\cite{thevenaz95:mackey}.
Besides the definition of Mackey functors sketched at the beginning of \cref{sec:mky-2-grps},
there is a more succinct formulation due to Lindner~\cite{lindner76:mackey}.

\begin{definition}\label{def:Mackey-functors}
  Let \(\Bbbk\) be a field and \(G\) a finite group.
  The category \(\mky_{\Bbbk}(G)\) of \emph{Mackey functors} of \(G\)
  is given by the \(\Bbbk\)\hyp{}linear functor category \([\Bbbk\Spiso{G},\kVect]\).
\end{definition}

Examples of Mackey functors are plentiful:
any representation of \(G\) defines one, see~\cite[Example~53.1]{thevenaz1995:G-algebrasModularRepresentation} as well as~\cite[Proposition~10.1]{panchadcharam07:mackey}.
For an extensive overview, we refer the reader to~\cite[Chapter~53]{thevenaz1995:G-algebrasModularRepresentation}.

\begin{example}\label{ex:mackey-algebra}
  To any finite group \(G\),
  one can associate a finite\hyp{}dimensional algebra \(\mathbb{M}G\)%
  —the \emph{Mackey algebra of \(G\)}—%
  whose category of modules is equivalent to \(\mky_{\Bbbk}(G)\), see~\cite[Propositions~3.1~and~3.2]{thevenaz95:mackey}.
  Its finite\hyp{}dimensional modules are in correspondence with the (pointwise) finite\hyp{}dimensional Mackey functors, which we denote by \(\mkyfin_{\Bbbk}(G)\).
\end{example}

Given that \(\Bbbk\Spiso{G}\) is symmetric monoidal, \(\mky_{\Bbbk}(G)\) is closed symmetric monoidal when equipped with Day convolution as its tensor product.\note{%
  As such, we refrain from adding ``left'' or ``right'' prefixes to the different duality notions when talking about Mackey functors.%
}
Furthermore, \(\Bbbk\Spiso{G}\) has a finite dense subcategory   whose objects form a complete set of representatives of transitive \(G\)\hyp{}sets.
The arguments of \cref{ex:mackey-dense-subcategories} may now be used to show that
the tensor product and internal hom of finite\hyp{}dimensional Mackey functors is finite\hyp{}dimensional,
see~\cite[Section~9]{panchadcharam07:mackey}.

The following result gives an alternative proof for the sketched argument
in~\cite[Lemma~2.2]{bouc05:mackey}
that the classes of rigid and finitely\hyp{}generated projective Mackey functors coincide.

\begin{proposition}\label{prop:mky-GV-and-rigid}
  Let \(G\) be a finite group and suppose \(\Bbbk\) is a field.
  The category \(\mkyfin_\Bbbk(G)\) is a \GV{} category with \(\rd*{\mkyfin_\Bbbk( \blank,1)}\) as dualising object.
  Its rigid objects are precisely the finitely\hyp{}generated projective Mackey functors.

  Furthermore, \(\mkyfin_{\Bbbk}(G)\) is rigid itself if and only if \(\mathbb{M}G\) is semisimple,
  which is equivalent to \(\mathrm{char}\,\Bbbk\) not dividing the order of \(G\).
\end{proposition}
\begin{proof}
  The fact that \(\mkyfin_{\Bbbk}(G)\) is a \GV{} category was proven in~\cite[Theorem~9.2]{panchadcharam07:mackey}.
  Alternatively, we may use that \(\Bbbk\Spiso{G}\) is a rigid category,
  see~\cite[Section~2]{panchadcharam07:mackey},
  and apply
  \cref{prop:fun-cat-GV} to obtain that \(\mkyfin_{\Bbbk}(G)\)
  is a \GV{} category with \(\rd*{\mkyfin_\Bbbk(\blank,1)}\in \mkyfin_{\Bbbk}(G)\) as its dualising object.
  Due to \cref{prop:characterisation-of-rigidity},
  a Mackey functor is rigid dualisable if and only if it is finitely\hyp{}generated projective.

  Now let us assume \(\mkyfin_{\Bbbk}(G)\) is rigid and recall that it is equivalent to the category \(\lmod{\mathbb{M}G}\) of finite\hyp{}dimensional modules of the Mackey algebra.
  Since \(\mathbb{M}G\) is finite\hyp{}dimensional and every object in \(\mkyfin_{\Bbbk}G\) is projective by \cref{prop:characterisation-of-rigidity}, any submodule of \(\mathbb{M}G\) must have a complement.
  In particular, \(\mathbb{M}G\) is semisimple.
  Conversely, in case \(\mathbb{M}G\) is semisimple, all objects of \(\lmod{\mathbb{M}G}\), and therefore also \(\mkyfin_{\Bbbk}(G)\), are projective.
  As they are furthermore finitely\hyp{}generated, \cref{prop:characterisation-of-rigidity} implies the rigidity of \(\mkyfin_{\Bbbk}(G)\).

  Corollary~14.4 of~\cite{thevenaz95:mackey}, shows that \(\mathbb{M}G\) is semisimple if \(\mathrm{char}\,\Bbbk\) does not divide \(|G|\).
  On the other hand, the category \(\mathrm{Rep}_{\Bbbk}(G)\) of finite\hyp{}dimensional representations of \(G\) over \(\Bbbk\) is a full subcategory of \(\mkyfin_{\Bbbk}(G)\), see~\cite[Proposition~10.1]{panchadcharam07:mackey}.
  Thus, if \(\mkyfin_{\Bbbk}(G)\) is semisimple, so is \(\mathrm{Rep}_{\Bbbk}(G)\),
  implying that \(\mathrm{char}\,\Bbbk\) does not divide \(|G|\)
  by Maschke's theorem.
\end{proof}

\subsection{Crossed modules}\label{sec:2-groups}

\textsc{Another example for \GV{} structures} on abelian \(\Bbbk\)\hyp{}linear functor categories arises from studying (strict) 2-groups,
which can be identified with crossed modules.
Our exposition follows~\cite{wagemann2021:CrossedModules}.

\begin{definition}\label{def:str-2-grp}
  A \emph{strict \(2\)\hyp{}group} is a (small) groupoid \(\cat{G}\) endowed with a strict monoidal structure,
  such that the monoid of objects \((\Ob(\cat{G}), \otimes,1)\) is a group.
\end{definition}

Every strict 2-group is a rigid category;
the left and right dual of an object \(g \in \cat{G}\) is given by its inverse \(g^{-1} \in \cat{G}\).

Let \(G\) be the group of objects of \(\cat{G}\) and \(H = \cat{G}(1, \blank)\) the set of arrows which start at the monoidal unit \(1\in G\).
Writing \(t \from H \to G \) for the \emph{target map},
we define a group structure on \(H\) via the multiplication
\[
  h' h \defeq (h' \otimes \id_{t(h)}) \circ h, \qquad \text{for all }h,h' \in H.
\]
For any \(g\in G\) and \(h\in H\) there is a unique \(h' \in H\) such that \(\id_g\otimes h = h' \otimes \id_g\).
This induces a map \(\alpha \from G \to \mathrm{Aut}(H)\) and turns the quadruple \((G, H, t, \alpha)\) into a crossed module as defined below.

\begin{definition}\label{def:crossed-module}
  A \emph{crossed module} is a quadruple \((G,H,t,\alpha)\)
  consisting of two groups \(G,H\)
  and two group homomorphisms \(t \from H \to G\), \(\alpha \from G \to \mathrm{Aut}(H)\),
  such that for all \(g \in G\) and \(h,l \in H\), we have
  \begin{equation}\label{eq:cross-mod-def-rel}
    t(\alpha(g)h) = g \circ t(h) \circ g^{-1}, \qquad
    \alpha(t(l))h = l \circ h \circ l^{-1}.
  \end{equation}
\end{definition}

\begin{remark}\label{rmk:crossed-modules-and-2-groups}
  Any crossed module \((G,H,t,\alpha)\) defines a strict \(2\)\hyp{}group \(\cat{G}\),
  with \(\Ob\cat{G}=G\), and \(\cat{G}(g,g') =\{\, (h,g) \in H \times G \mid t(h)g=g'\,\}\).
  Composition is given by \((h', t(h)g)\circ(h,g) \defeq (h'h,g)\),
  and for the tensor product we have \((h,g) \otimes (h', g') \defeq (h\,\alpha(g)h', g g')\).
  The left and right dualising functors coincide;
  they map any object \(b \in G\) to \(\ld{b}=b^{-1}=\rd{b}\).
  For a morphism \(f=(h,b)\from b \to c\), we have
  \[
    \ld{\!{}f}  =(\alpha(c^{-1})h, c^{-1}) = (\alpha(b^{-1})h, c^{-1}) = \rd{f}.
  \]
\end{remark}

There is an equivalence of categories between crossed modules and strict 2-groups as is shown for example in~\cite[Theorem~1.9.2]{wagemann2021:CrossedModules}.

\begin{example}\label{ex:2-groups-in-algebra-and-rep-theo}\hfill
  \begin{itemize}
    \item\emph{Hopf--Galois extensions and skew braces}:
    the \emph{holomorph} of a group \(H\) is the semidirect product \(H \rtimes \mathrm{Aut}(H)\).
    The two group homomorphisms \(t \from H \to \mathrm{Aut}(H)\), \(t(l)h = lhl^{-1}\), and \(\alpha=\id\from \mathrm{Aut}(H) \to \mathrm{Aut}(H)\)
    turn \((\mathrm{Aut}(H), H, t,\id_{\mathrm{Aut}(H)})\) into a crossed module.
    As explained in the Introduction of~\cite{byott2024:insolubleHolomorph}
    and~\cite[Section~2]{bachiller2016:CounterexampleConjectureBraces},
    holomorphs can be used to classify certain Hopf--Galois extensions as well as skew-braces.

    \item\emph{Representations of finite groups}:
    a short exact sequence of groups
    \[
      \begin{mytikzcd}[ampersand replacement=\&]
        0 \& C \& E \& G \& 0
        \arrow[from=1-1, to=1-2]
        \arrow["\iota", from=1-2, to=1-3]
        \arrow["t", from=1-3, to=1-4]
        \arrow[from=1-4, to=1-5]
      \end{mytikzcd}
    \]
    such that \(C\) embeds into the centre of \(E\) is called a \emph{central extension} of \(G\).
    Let \(\kappa \from G\to E\) be a set-theoretical section of \(t\from E\to G\).
    A direct computation shows that the map \(\alpha\from G \to \mathrm{Aut}(E)\)
    defined by \(\alpha(g)e \defeq \kappa(g)e{\kappa(g)}^{-1}\) is independent of the choice of the section
    and a homomorphism of groups.
    The tuple \((E,G,t,\alpha)\) forms a crossed module, see~\cite[Section 1.3]{wagemann2021:CrossedModules}.

    As discussed in~\cite[Chapter~1]{hoffman-humphreys1992:ProjectiveRepresentationsSymmetricGroup},
    if \(G\) is a finite group, then
    any projective representation \(\rho \from G \to \mathrm{PGL}_{\complex}(V)\) of \(G\)
    can be lifted to a linear representation \(\varrho \from E \to \GL(V)\) of a certain central extension \(E\) of \(G\).
    Projective representations themselves are studied in the context of representation theory of semidirect products,
    see for example~\cite{ceccherini-silberstein-scarabotti-tolli2022:RepresentationFiniteGroupExtensions}.
  \end{itemize}
\end{example}

Let \((G,H,t,\alpha)\) be a crossed module whose associated \(2\)\hyp{}group we denote by \(\cat{G}\).
Using the results of \cref{sec:tensor-rep-functor-cats},
we determine a \GV{} structure on the category
\(\langle\cat{G}, \kvect\rangle\) of ordinary functors between \(\cat{G}\) and \(\kvect\).

Hereto, we need to analyse the structure of \(\cat{G}\) in more detail.
By \cref{eq:cross-mod-def-rel}, the image \(K \defeq \im t \subset G\) is a normal subgroup of \(G\)
and the kernel \(L \defeq \ker t\) is normal and central in \(H\).
The latter follows from the identity
\[
  lhl^{-1} =\alpha(t(l))h = \alpha(e)h = h \qquad\quad \text{ for all } l\in L, h\in H.
\]
For any \(k=t(h) \in K\) and \(l\in L\), we get \(\alpha(k)l = \alpha(t(h))l = hlh^{-1}= l\).
Hence,
there is a unique group homomorphism
\(\overline{\alpha} \from Q \defeq G/K \to \mathrm{Aut}(L)\), with
\[
  \overline{\alpha}([g])l=\alpha(g)l \in L \subset H,
\]
for all \(g\in G\) and \(l\in L\).

The connected components of \(\cat{G}\) are in bijection with the elements of \(Q\).
Given any element \(g\in G\), we write \( \cat{G}_{g}\) for the maximal full connected pointed subgroupoid of \(\cat{G}\) whose distinguished object is \(g\).
The set of objects of \(\cat{G}_{g}\) is \(Kg\) and its morphisms correspond to \(H \times Kg\).

\begin{lemma}\label{lem:canonical-inclusion}
  Let \(g \in G\) and write \(\mathbf{B}L\) for the delooping of \(L\) with single object \(g\).
  The canonical inclusion \(\Theta_g \from \mathbf{B}L \emb \cat{G}_g\) sending \(g\) to \(g\) and \(l \from g\to g\) to \((l,g) \from g\to g\) is
  essentially surjective and fully faithful—an equivalence.
\end{lemma}

To specify a quasi-inverse,
fix a set-theoretical section \(\iota \from K \to H\) of
the surjective homomorphism of groups \(t\from H \to K\).
Define the quasi-inverse \(\Psi_g \from \cat{G}_g\to \mathbf{B}L\)
by mapping each object to \(g\) and any morphism \((h,kg)\) to \(\iota({(t(h)k)}^{-1}) h \iota(k)\).
Using that \(\langle \mathbf{B}L, \kvect\rangle\) can be identified with the category \(\mathsf{rep}(L)\) of finite\hyp{}dimensional representations of \(L\),
the pushforwards of \(\Theta_g\) and \(\Psi_g\) establish an equivalence of categories
\begin{equation}\label{eq:toopushy}
  \begin{aligned}
    \Theta_{g}^{*} \colon \langle \cat{G}_g, \kvect\rangle &\ \rightleftarrows\ \mathsf{rep}(L) \cocolon \Psi_g^{*} \\
    F &\mapsto (Fg,\rho_{F}) \quad \text{(where \(\rho_F(l) = F(l,g)\))} \\
    \left\{\begin{aligned}
      kg &\mapsto M\\
      (hk,g) &\mapsto \rho\left(\iota({(t(h)k)}^{-1}) h \iota(k)\right)
    \end{aligned}\right\} & \longmapsfrom (M, \rho).
  \end{aligned}
\end{equation}

\begin{notation}
  Define \(\mathsf{rep}_{\cat{G}}(L) \defeq \bigoplus_{q\in Q}{\mathsf{rep}(L)}_{q}\),
  where  \({\mathsf{rep}(L)}_{q} = \mathsf{rep}(L)\).
  Given some \(M\in \mathsf{rep}_{\cat{G}}(L)\), we write \(M_q\) for its homogeneous component of degree \(q\in Q\).
  Note that in case \(K\) has finite index in \(G\), we have
  \[
    \langle \cat{G}, \kvect\rangle \cong \bigoplus_{[g]\in Q}\langle \cat{G}_g, \kvect \rangle.
  \]
\end{notation}

\begin{lemma}\label{lem:structure-of-G-funs}
  Let \(\cat{G}\) be a strict \(2\)\hyp{}group, and write \((G, H, t, \alpha)\) for the corresponding crossed module.
  If \(K=\im t\subset G\) has finite-index,
  then \(\mathsf{rep}_{\cat{G}}(L) \simeq \langle \cat{G}, \kvect\rangle\).
\end{lemma}

The tools developed in \cref{sec:tensor-rep-functor-cats} allow us to endow \(\langle\cat{G}, \kvect\rangle\)
with the structure of a \GV{} category which we will transfer to \(\mathsf{rep}_{\cat{G}}(L)\).
The tensor product obtained in this manner will permute the homogeneous components and ``twist'' the action of \(L\) by virtue of \(\overline{\alpha}\).
In this regard, we introduce the following notation:
for any \(q \in Q\) and \(M \in \mathsf{rep}(L)\),
we denote by \(M^{\overline{\alpha}(q)}\) the representation of \(L\) whose underlying vector space is \(M\), endowed
with the action \(l \blact m = \overline{\alpha}(q)l \lact m\) for all \(l \in L\) and \(m \in M\).

\begin{proposition}\label{prop:2-groups-tensor-product}
  Suppose \((G,H,t,\alpha)\) is a crossed module with \(G\) and \(H\) finite and let \(\cat{G}\) be its associated strict \(2\)\hyp{}group.
  The category \(\mathsf{rep}_{\cat{G}}(L)\) is a right \(r\)\hyp{}category.
  Its tensor product and dualising functor are defined by the assignments
  \[
    M\otimes N \defeq (M\otimes_{\Bbbk L} N^{\overline{\alpha}(p^{-1})})\in {\mathsf{rep}_{\cat{G}}(L)}_{pq}, \qquad
    D M \defeq \rd*{\big(M^{\overline{\alpha}(p)}\big)} \in {\mathsf{rep}_{\cat{G}}(L)}_{p^{-1}}
  \]
  for \(p, q \in Q\) and \(M \in {\mathsf{rep}_{\cat{G}}(L)}_{p}\), \(N \in {\mathsf{rep}_{\cat{G}}(L)}_{q}\).
\end{proposition}
\begin{proof}
  By \cref{prop:fun-cat-GV},
  \(\langle\cat{G},\kvect\rangle \cong [\Bbbk\cat{G}, \kvect]\) can be endowed with the structure of a right \GV{} category
  with \(\rd*{\Bbbk\cat{G}( \blank , 1)}\) as dualising object;
  write \(R \from \langle\cat{G},\kvect\rangle^{\op} \to \langle\cat{G},\kvect\rangle\) for its dualising functor.

  By \cref{lem:structure-of-G-funs}, there are \(\Bbbk\)\hyp{}linear equivalences
  \[
    \adj{\widehat{\Theta}}{\widehat{\Psi}}{[\Bbbk\cat{G}, \kvect]}{\mathsf{rep}_{\cat{G}}(L)}
  \]
  that are determined
  by the pushforwards of the equivalences
  \[
    \adj{\Theta_g}{\Psi_{g}}{\langle\cat{G}_g, \kvect\rangle}{\mathbf{B}L},
  \]
  for each homogeneous component \([\Bbbk\cat{G}_{[g]}, \kvect]\).
  Using the formulas of \cref{eq:toopushy},
  we can therefore explicitly transfer the \GV{} structure of \([\Bbbk\cat{G}, \kvect]\) to \(\mathsf{rep}_{\cat{G}}(L)\).

  The dualising object is mapped to
  \(\Theta(\rd*{\Bbbk\cat{G}( \blank, 1)})=\rd*{\Bbbk L} \cong \Bbbk L\in {\mathsf{rep}_{\cat{G}}(L)}_{[1]}\),
  where we used that  \(\Bbbk L\) is a Frobenius algebra for the last equality.
  The tensor product and right dualising functor of \(\mathsf{rep}_{\cat{G}}(L)\) are given by
  \begin{align*}
    \otimes &\defeq \widehat{\Theta} * (\widehat \Psi \times \widehat \Psi) \from \mathsf{rep}_{\cat{G}}(L) \times \mathsf{rep}_{\cat{G}}(L) \to \mathsf{rep}_{\cat{G}}(L), \\
    D &\defeq \widehat{\Theta}\mathsf{R}\widehat \Psi^{\op} \from {\mathsf{rep}_{\cat{G}}(L)}^{\op} \to \mathsf{rep}_{\cat{G}}(L).
  \end{align*}
  In order to compute it explicitly, we consider two elements \(p,q\in Q\) as well as two representations \((M, \rho)\in {\mathsf{rep}_{\cat{G}}(L)}_{p}\) and \((N, \tau) \in {\mathsf{rep}_{\cat{G}}(L)}_{q}\).
  The Day convolution of the functors corresponding to \(M\) and \(N\) is computed via a certain colimit and \(\widehat{\Theta}\), as an equivalence of categories, commutes with this colimit.
  Thus, by \cref{eq:day-convolution-simp-right} and the definition of coends,
  the homogeneous component \({(M\otimes N)}_x\) of degree \(x\in Q\) is the coequaliser of
  \[
    \begin{mytikzcd}[ampersand replacement=\&]
      \displaystyle{\coprod_{r, s \in Q}}  \Bbbk L_{r,s} \otimes_{\Bbbk} M_r \otimes_{\Bbbk} N_{s^{-1}x} \& \displaystyle{\coprod_{r\in Q}  M_r \otimes_{\Bbbk} N_{r^{-1}x},}
      \arrow[shift right=1, from=1-1, to=1-2]
      \arrow[shift left=1, from=1-1, to=1-2]
    \end{mytikzcd}
  \]
  where \(L_{r,s}=L\) if \(r=s\) and \(\varnothing\) otherwise.
  Its parallel morphisms are
  \[
    \begin{aligned}
      l\otimes_{\Bbbk} m\otimes_{\Bbbk} n &\mapsto  F\rho(l)m \otimes_{\Bbbk} n, \\
      l\otimes_{\Bbbk} m\otimes_{\Bbbk} n &\mapsto  m \otimes \tau(\rd{l})n = m \otimes \tau(\overline{\alpha}(s^{-1})l)n,
    \end{aligned}
  \]
  where \(l \in L_{r,s}\), \(m \in M_{r}\) and \(n \in N_{s^{-1}x}\).

  In order to compute \(D\), we note that \(\mathsf{R}(F)= \rd*{F(\ld{\blank})}\) for any \(F \in [\Bbbk\cat{G}, \kvect]\).
  Thus, given \(x\in Q\), we compute
  \[
    {DM}_x =
    \begin{cases}
      \rd*{M} & x^{-1} = p, \\
      \{0\} & \text{otherwise}.
    \end{cases}
  \]
  Hence, the action of any \(l \in L\) on \(\rd*{M}\) is given by \(\rd*{\big(\rho(\overline{\alpha}(p)l)\big)}\from \rd*{M}_{p} \to \rd*{M}_{p}\).
\end{proof}

\begin{proposition}\label{prop:2-grp-GV-and-rigid}
  Let \((G,H,\alpha,t)\) be a crossed module with \(G\) and \(H\) finite and write \( \cat{G}\) for its associated strict 2-group.
  The category \(\mathsf{rep}_{\cat{G}}(L)\) is rigid if and only if \(\mathrm{char}\, \Bbbk\) does not divide the order of \(L= \ker t\).
\end{proposition}
\begin{proof}
  It follows from \cref{prop:characterisation-of-rigidity} that \(\mathsf{rep}_{\cat{G}}(L)\) is rigid
  if and only if all of its objects are finitely\hyp{}generated and projective.
  As \(\mathsf{rep}_{\cat{G}}(L)\) is a direct sum of finitely many copies of \(\mathsf{rep}(L)\)
  and every object of \(\mathsf{rep}(L)\) is finitely\hyp{}generated,
  this is the case if and only if all objects of \(\mathsf{rep}(L)\) are projective.
  The latter is equivalent to \(\Bbbk L\) being semisimple, which
  corresponds to \(\mathrm{char}\, \Bbbk\) and \(|L|\) being coprime
  by Maschke's theorem.
\end{proof}

%%% Local Variables:
%%% mode: LaTeX
%%% TeX-master: "../main"
%%% TeX-command-extra-options: "-shell-escape -fmt=prec.fmt -file-line-error -synctex=1"
%%% End:
